\section{Gauss 和与 Jacobi 和}\label{VI:4}
\origpage{202}{196}

\subsection{}\label{VI:4-1}
设 $k$ 是特征为 $p$ 的有限域，$\chi$ 是 $k^\times$ 的一个特征，$\psi$ 是 $k$ 的加法群的一个非平凡特征。我们采用如下\term{Gauss 和}{Gauss sum}{somme de Gauss}的定义：
\begin{equation*}\tag{4.1.1}\label{VI:eq:4-1-1}
  \tau(\chi,\psi)=-\sum_{x\in k^\times}\psi(x)\chi^{-1}(x).
\end{equation*}
请注意这里的负号。通常取 $\chi$ 和 $\psi$ 为复值特征；我们取它们的值于 $E_\lambda^\times$（参见 \ref{VI:1-15} 的证明）。把 $-\tau(\chi,\psi)$ 看成对 $k$ 上概形 $\dG_m$ 的有理点求和。由第 \ref{VI:1} 节，
$-\tau(\chi,\psi)=\tr(F^\ast,\h_c^\ast(\dG_m,\sF(\psi\chi^{-1})))$。此外还有如下结论。

\begin{proposition_}\label{VI:4-2}
$\dG_m$ 上以 $\sF(\psi\chi^{-1})$ 为系数的上同调满足：
\begin{enumerate}[(i)]
\item 当 $i\ne1$ 时，$\h_c^i=0$，且 $\dim\h_c^1=1$；
\item $F^\ast$ 在 $\h_c^1$ 上的作用是乘以 $\tau(\chi,\psi)$；
\item 若 $\chi$ 非平凡，则 $\h_c^\bullet\iso\h^\bullet$。
\end{enumerate}
\end{proposition_}
\begin{proof}
对每个与 $p$ 互素的 $n$，记 $K_n$ 为 Kummer 正合列
$0\to\dmu_n\to\dG_m\xrightarrow{x^n}\dG_m\to0$
所定义的 $\dG_m$ 上的 $\dmu_n$-扭子。若 $k$ 有 $q$ 个元素，则在 $k$ 上有 $\dmu_{q-1}=k^\times$，而 $\dG_m/k$ 上的 Lang 扭子就是 $K_{q-1}$。因此 $\sF(\chi^{-1})=\chi(K_{q-1})$。断言 \ref{VI:4-2}(ii) 可由 (i)、(iii) 得到，后两者则包含在下面更具几何性的命题中。其中，$\sF(\psi)$ 表示将 $\dG_m/k$ 上的 $\sF(\psi)$ 从 $k$ 扩张标量到 $\dF$ 所得的 $\dG_m/\dF$ 上的层（\ref{VI:1-8}.b）。
\end{proof}

\begin{proposition_}\label{VI:4-3}
设 $\chi:\dmu_n\to E_\lambda^\times$。$\dG_m$ 上以 $\sF(\psi)\otimes\chi(K_n)$ 为系数的上同调满足：
\begin{enumerate}[\indent (i)]
\item 当 $i\ne1$ 时，$\h_c^i=0$，且 $\dim\h_c^1=1$；
\item 若 $\chi$ 非平凡，则 $\h_c^\bullet\iso\h^\bullet$。
\end{enumerate}
\end{proposition_}
\begin{proof}
层 $\sF(\psi)$ 是 $\dG_a$ 上一个局部常值层在 $\dG_m$ 上的限制；这个层在无穷远处野分歧，Swan 导子为 $1$。若 $\chi=1$，则 $\chi(K_n)$ 为常值层；若 $\chi\ne1$，则它在 $0$ 和 $\infty$ 处均为温和分歧。

因此，$\sF(\psi)\otimes\chi(K_n)$ 在 $\infty$ 处分歧。应用 \ref{VI:1-18}(b,c)，可得当 $i\ne1$ 时，$\h_c^i(\dG_m,\sF(\psi)\otimes\chi(K_n))=0$。若 $\chi\ne1$，则此层在 $0$ 和 $\infty$ 处都分歧，由 \ref{VI:1-19-1} 即得 (ii)。

它在 $0$ 和 $\infty$ 处的 Swan 导子分别为 $0$ 和 $1$。由 \eqref{VI:eq:3-2-1}，Euler--Poincaré 示性数为 $-1$，于是命题得证。
\end{proof}

\subsection{备注}\label{VI:4-4}
若 $\chi$ 非平凡，\ref{VI:4-2}(iii) 与 Poincaré 对偶表明，
$\h_c^1(\dG_m,\sF(\psi\chi^{-1}))$ 和
$\h_c^1(\dG_m,\sF(\psi^{-1}\chi))$
互为对偶，对偶取值于 $E_\lambda(-1)$。因此
$\tau(\chi,\psi)\tau(\chi^{-1},\psi^{-1})=q$，即
$|\tau(\chi,\psi)|=q^{1/2}$。
\translatornote{原书此处将绝对值误印为 $q$。在复嵌入下，有限阶特征的值为单位根，故 $\overline{\tau(\chi,\psi)}=\tau(\chi^{-1},\psi^{-1})$。由前一个等式得到 $|\tau(\chi,\psi)|^2=q$，所以应为 $q^{1/2}$。例如 $k=\dF_3$，取二次特征与 $\psi(x)=\exp(2\pi i x/3)$，便得到绝对值为 $\sqrt3$ 的 Gauss 和。}

\subsection{}\label{VI:4-5}
若 $k$ 是 $\dF_q$ 的 $N$ 次扩张，也可将 \eqref{VI:eq:4-1-1} 看成 $\dF_q$ 上的一个 $N$ 元和。更一般地，设 $k$ 是 $\dF_q$ 上次数为 $N$ 的平展代数。它是若干域 $k_i$ 的直积，其中 $[k_i:\dF_q]=N_i$。令
\begin{equation*}\tag{4.5.1}\label{VI:eq:4-5-1}
\varepsilon(k)=(-1)^{\sum(N_i+1)}.
\end{equation*}
这就是 Frobenius 代换 $\varphi\in\gal(\dF/\dF_q)$ 在
$S=\hom_{\dF_q}(k,\dF)$ 上所诱导置换的符号。

设 $\psi:\dF_q\to E_\lambda^\times$ 是非平凡特征，$\chi:k^\times\to E_\lambda^\times$。令
\begin{equation*}\tag{4.5.2}\label{VI:eq:4-5-2}
\tau_{\dF_q}(\chi,\psi)=(-1)^N\sum_{x\in k^\times}
 \psi\bigl(\tr_{k/\dF_q}(x)\bigr)\chi^{-1}(x).
\end{equation*}
若 $\chi$ 的各坐标为 $\chi_i:k_i^\times\to E_\lambda^\times$，则显然有恒等式
\begin{equation*}\tag{4.5.3}\label{VI:eq:4-5-3}
\tau_{\dF_q}(\chi,\psi)=\varepsilon(k)\prod_i\tau(\chi_i,\psi\circ\tr_{k_i/\dF_q}).
\end{equation*}
在作一些准备之后，我们将在 \ref{VI:4-10} 给出和式 \eqref{VI:eq:4-5-2} 的上同调解释。在 \ref{VI:4-12}，我们将把 Hasse--Davenport 恒等式解释为 \eqref{VI:eq:4-5-3} 下述特例的一种扭曲形式（参见 \ref{VI:1-12}）：若 $k=\dF_q^N$，而 $\chi$ 是 $\dF_q^\times$ 的特征，则
$\tau_{\dF_q}(\chi\circ N_{k/\dF_q},\psi)=\tau(\chi,\psi)^N$。

\subsection{}\label{VI:4-6}
回顾如下构造。若 $M$ 是环 $A$ 上的有限生成投射模，则从 $\spec(A)$ 上的仿射概形到集合范畴 $\mathsf{Ens}$ 的函子
$\spec(B)\mapsto M\otimes_A B$
由仿射概形 $\dV(M^\vee)$ 表示（记号沿用 EGA），即
$\operatorname{Sym}_A^\bullet(M^\vee)$ 的谱。若 $M=A^n$，这就是标准的 $n$ 维仿射空间。

设 $M$ 是带单位的 $A$-代数，并令 $V=\dV(M^\vee)$。按定义，对 $A$ 的每个扩张 $B$，$V$ 的 $B$ 值点集 $V(B)$ 就是 $M\otimes_A B$。由于函子 $B\mapsto M\otimes_A B$ 取值于带单位的环，$V$ 是 $\spec(A)$ 上的带单位环概形。范数和迹映射 $M\otimes_A B\to B$ 关于 $B$ 具有函子性，因此分别对应于概形态射 $N,T:V\to\dG_a$。我们需要考虑由 $V$ 得到的下列概形：
\begin{enumerate}[a)]
\item $V^\ast$ 是 $V$ 的可逆元所成的开子概形，它关于乘法 $\cdot$ 是一个群概形；
\item $W$ 是方程 $T=0$ 定义的超平面，且 $W^\ast=W\cap V^\ast$；
\item $P$ 是 $\spec(A)$ 上仿射空间 $V$ 的无穷远超平面，即 $(V\setminus\{0\})/\dG_m$。若 $Q$ 是 $W$ 的无穷远超平面，则 $W^\ast/\dG_m$ 和 $V^\ast/\dG_m$ 分别是 $\spec(A)$ 上射影空间 $Q$ 和 $P$ 的开子概形。
\begin{equation*}\tag{4.6.1}\label{VI:eq:4-6-1}
\xymatrix{
 W^\ast\ar@{^{(}->}[r]\ar[d]^-\pi&V^\ast\ar[d]^-\pi\\
 W^\ast/\dG_m\ar@{^{(}->}[r]&V^\ast/\dG_m
}\qquad\subset\qquad
\xymatrix{
 W\setminus\{0\}\ar@{^{(}->}[r]\ar[d]^-\pi&V\setminus\{0\}\ar[d]^-\pi\\
 Q\ar@{^{(}->}[r]&P.
}
\end{equation*}
\end{enumerate}
若 $M=A^I$，其中 $I$ 为有限集，则 $V\simeq\dG_a^I$，$V^\ast=\dG_m^I$；$V^\ast/\dG_m$ 是环面 $\dG_m^I/(\text{对角嵌入的 }\dG_m)$，$N$ 和 $T$ 分别写成 $\prod x_i$ 和 $\sum x_i$，故 $Q$ 是 $P$ 中由 $\sum x_i=0$ 定义的射影超平面。

我们关心的是 $M$ 在 $A$ 上有限平展的情形。此时上述描述在 $\spec(A)$ 的平展拓扑下局部成立，因此可以写 $V^\ast=\dG_m^I$，其中 $I$ 是 $\spec(A)$ 上的有限集的局部常值层。例如，若 $A$ 是域，$\bar A$ 是其代数闭包，而 $M$ 是可分 $A$-代数，则令 $I=\hom_A(M,\bar A)$；在 $\bar A$ 上有 $V\simeq\dG_a^I$。在此同构下，$N$ 和 $T$ 分别为 $\prod x_i$ 和 $\sum x_i$。

\subsection{Kummer 扭子}\label{VI:4-7}
设 $S$ 是概形，$n$ 是在 $S$ 上可逆的整数，$G$ 是 $S$ 上纤维连通的交换群概形，采用乘法记号。列
\[\xymatrix{0\ar[r]&G_n\ar[r]&G\ar[r]^-{x^n}&G\ar[r]&0}\]
是正合的。它定义了 $G$ 上的一个 $G_n$-扭子 $K_n(G)$，也简记为 $K_n$。对于从 $S$ 上的平展层 $G_n$ 到常值层 $E_\lambda^\times$ 的同态 $\chi:G_n\to E_\lambda^\times$，用 $\sK_n(\chi)$ 表示 $E_\lambda$-层 $\chi^{-1}(K_n)$。

这些扭子 $K_n$ 构成群逆系统 $G_n$ 下的扭子逆系统，其过渡态射为 $x\mapsto x^d:G_{nd}\to G_n$。这表达为下图的交换性：
\[\xymatrix{
0\ar[r]&G_{nd}\ar[r]\ar[d]^-{x^d}&G\ar[r]^-{x^{nd}}\ar[d]^-{x^d}&G\ar[r]\ar@{=}[d]&0\\
0\ar[r]&G_n\ar[r]&G\ar[r]^-{x^n}&G\ar[r]&0.
}\]
因而 $\sK_n(\chi)=\sK_{nd}(\chi\circ x^d)$。

\subsection{}\label{VI:4-8}
将 \ref{VI:4-6} 的构造用于 $A=\dF_q$ 及 $\dF_q$ 上的平展代数 $M=k$。依照一般约定，在 \ref{VI:4-6} 中记作 $V,V^\ast,\ldots$ 的 $\dF_q$-概形现添下标 $0$；相同字母不带下标时，表示将其扩张标量到 $\dF$ 所得的概形。

令 $I=\hom(k,\dF)$，于是 $V\simeq\dG_a^I$，且 $T$ 为 $(x_i)\mapsto\sum x_i$。因此
\begin{equation*}\tag{4.8.1}\label{VI:eq:4-8-1}
\sF(\psi\circ\tr_{k/\dF_q})=T^\ast\sF(\psi)=\bigotimes_i\pr_i^\ast\sF(\psi).
\end{equation*}
由于 $V^\ast\simeq\dG_m^I$，$V_n^\ast\simeq\dmu_n^I$ 的一个 $E_\lambda^\times$ 值特征 $\chi$，可以写成由 $I$ 编号的一族 $\dmu_n$ 的特征 $(\chi_i)_{i\in I}$。它在 $\dF_q$ 上有定义，当且仅当对于每个 $\sigma\in\gal(\dF/\dF_q)$，有 $\chi_{\sigma i}=\chi_i\circ\sigma^{-1}$。此时，它定义了 $V_0^\ast$ 上的一个 $E_\lambda$-层 $\sK_n((\chi_i)_{i\in I})$。若各 $\chi_i$ 的乘积平凡，则 $\chi$ 通过 $(V_0^\ast/\dG_m)_n$ 的一个特征分解，而 $\sK_n((\chi_i)_{i\in I})$ 是 $V_0^\ast/\dG_m$ 上某层的逆像；后者仍记作同一符号。此构造与 Lang 扭子之间有如下关系。

\begin{lemma_}\label{VI:4-9}
\begin{enumerate}[(i)]
\item 当 $n$ 具有充分多的因子时，有交换图
\[\xymatrix{
0\ar[r]&(V_0^\ast)_n\ar[r]\ar[d]^-\tau&V_0^\ast\ar[r]^-{x^n}\ar[d]^-\tau&V_0^\ast\ar[r]\ar@{=}[d]&0\\
0\ar[r]&k^\times\ar[r]&V_0^\ast\ar[r]^-\fL&V_0^\ast\ar[r]&0.
}\]
若 $\chi:k^\times\to E_\lambda^\times$ 是特征，则 $\sF(\chi)=\sK_n(\chi\circ\tau)$。
\item 当 $n$ 具有充分多的因子时，$\tau$ 将 $k^\times$ 等同于 $\gal(\dF/\dF_q)$ 作用在 $V_n^\ast\simeq\dmu_n^I$ 上的\term{余不变量}{coinvariants}{coinvariants}。
\item 若 $k$ 为域，$N=[k:\dF_q]$，$\omega\in I$ 是 $k$ 到 $\dF$ 的一个嵌入，且 $n=q^N-1$，则 $\chi\circ\tau$ 的 $\omega$ 分量为 $\chi\circ\omega^{-1}$。
\item 若 $\chi$ 在 $k$ 的每个因子上均非平凡，则所有 $(\chi\circ\tau)_i$ 都非平凡。
\end{enumerate}
\end{lemma_}
只需在 $k$ 是 $\dF_q$ 的 $N$ 次域扩张时证明引理。此时，$n$ 具有“充分多的因子”意指 $q^N-1\mid n$。选取嵌入 $\omega:k\hookrightarrow\dF$，便可将 $I$ 等同于 $\dZ/N$：$i\in\dZ/N$ 对应于 $\omega_i=\omega^{q^i}$。经同构 $V_0^\ast(\dF)=(\dF^\times)^I$，有
\begin{enumerate}[\indent a)]
\item $x\in k^\times$ 对应于 $(\omega_i(x))\in(\dF^\times)^I$；
\item $F((x_i)_{i\in\dZ/N})=(x_{i-1}^q)_{i\in\dZ/N}$。
\end{enumerate}
$V_n^\ast\simeq\dmu_n^I$ 的特征 $\chi=(\chi_i)$ 在 $\dF_q$ 上有定义，当且仅当 $\chi_{-i}(x)=\chi_0(x^{q^i})$（$i\in\dZ$）。若 $q^N-1\mid n$，这样的特征共有 $q^N-1$ 个：$\chi_0$ 通过 $\dmu_{q^N-1}$ 分解，并决定所有 $\chi_i$。因此，要证明 (ii)，只需证明 (i) 和 (iii)，或用其推论 (iv)；后者保证 $\chi\mapsto\chi\circ\tau$ 为单射。

现对 $n=q^N-1$ 证明 (i)、(iii)。将 $V^\ast\simeq\dG_m^I$ 的自身态射群用加法记号表示，特别将算子 $x\mapsto x^n$ 记为 $n$。若 $\alpha$ 是循环置换算子 $(x_i)\mapsto(x_{i-1})$，则
$q^N-1=(q\alpha)^N-1=(q\alpha-1)((q\alpha)^{N-1}+\cdots+1)$。
这就决定了 $\tau$：其 $i$ 分量为
$\tau((x_i))_i=\prod_{0\le j<N}x_{i-j}^{q^j}$，而由它诱导的 $(V_0^\ast)_n\to k^\times$ 写成
$\omega(x)=\prod_i x_{-i}^{q^i}$（参见上面的 a)）；由此即得 (iii)。

\begin{proposition_}\label{VI:4-10}
$V^\ast$ 上以 $\sF(\chi^{-1}\cdot\psi\tr_{k/\dF_q})$ 为系数的上同调满足：
\begin{enumerate}[\indent (i)]
\item 当 $i\ne N$ 时，$\h_c^i=0$，且 $\dim\h_c^N=1$；
\item 在 $\h_c^N$ 上，$F^\ast$ 的作用是乘以 $\tau_{\dF_q}(\chi,\psi)$；
\item 若 $\chi$ 在 $k$ 的每个因子上均非平凡，则 $\h_c^\bullet\iso\h^\bullet$。
\end{enumerate}
\end{proposition_}
应用 \ref{VI:4-9}(i)：在 $\dF$ 上，若 $\chi\circ\tau=(\chi_i)_{i\in I}$，则 $\sF(\chi)=\sK_n((\chi_i)_{i\in I})$。因此，(i)、(iii) 由下面更具几何性的命题推出（对于 (iii)，使用 \ref{VI:4-9}(iv)），而 (ii) 再由迹公式得到。

\begin{proposition_}\label{VI:4-11}
设 $(\chi_i)_{i\in I}$ 是 $\dmu_n(\dF)$ 的一族特征。$V^\ast\simeq\dG_m^I$ 上以
$\sK_n((\chi_i)_{i\in I})\otimes\sF(\psi\tr_{k/\dF_q})$
为系数的上同调满足：
\begin{enumerate}[\indent (i)]
\item 当 $i\ne N$ 时，$\h_c^i=0$，且 $\dim\h_c^N=1$；
\item 若所有 $\chi_i$ 均非平凡，则 $\h_c^\bullet\iso\h^\bullet$。
\end{enumerate}
\end{proposition_}
我们有 $V^\ast\simeq\dG_m^I$，
$\sK_n((\chi_i))=\bigotimes_i\pr_i^\ast\chi_i(K_n(\dG_m))$，以及
$\sF(\psi\circ\tr_{k/\dF_q})=T^\ast\sF(\psi)=\bigotimes_i\pr_i^\ast\sF(\psi)$。
于是可以应用 Künneth 公式，\ref{VI:4-11} 便由 \ref{VI:4-3} 得到。

\subsection{}\label{VI:4-12}
由 \ref{VI:2-4*} 及（\hyperref[IV]{Cycle}，\ref{IV:1-3} 例 2）还可推出：$I$ 中满足 $\chi_i=\chi_{\sigma i}$（$i\in I$）的置换 $\sigma$ 所成的群，在 $\h_c^N$ 上通过乘以符号 $\varepsilon(\sigma)$ 作用。我们由此给出 Hasse--Davenport 恒等式的另一证明。

若 $\chi$ 是 $\dF_q^\times$ 的特征，由于在 $\dF$ 上 $N$ 写成 $(x_i)\mapsto\prod x_i$，故
$\sF(\chi\circ N)=N^\ast\sF(\chi)=\bigotimes_i\pr_i^\ast\sF(\chi)$。
Künneth 公式给出
\[
\h_c^\bullet(V^\ast,\sF(\chi^{-1}\circ N\cdot\psi\circ\tr))
\simeq\h_c^\bullet(\dG_m,\sF(\chi^{-1}\psi))^{\otimes I},
\]
其中右侧的张量积按 \ref{VI:2-4*} 的意义理解。由于 $\h_c^1(\dG_m,\sF(\psi\chi^{-1}))$ 为一维，其通常意义下的 $\otimes I$ 张量幂只依赖 $I$ 的基数 $N$。应用（\hyperref[IV]{Cycle}，\ref{IV:1-3} 例 2），得到典范同构
\begin{equation*}\tag{4.12.1}\label{VI:eq:4-12-1}
\h_c^N\left(V^\ast,\sF(\chi^{-1}\circ N)\otimes T^\ast\sF(\psi)\right)
\simeq\h_c^1(\dG_m,\sF(\psi\chi^{-1}))^{\otimes N}
 \otimes_\dZ\bigwedge^N\dZ^I.
\end{equation*}
\translatornote{原书式 (4.12.1) 左侧将 $V^\ast$ 误印作 $V$，右侧漏印 $N$ 次张量幂。由前一行的分次 Künneth 同构，唯一非零项是 $N$ 个一次上同调的张量积，置换因子由行列式直线 $\bigwedge^N\dZ^I$ 记录，故应为这里的形式。下文 Frobenius 特征值的 $N$ 次幂也要求这一修正。原书同段将 $N^\ast\sF(\chi)$ 的各因子误写为 $\sF(\psi)$，亦已按拉回公式改正。}
为了计算 $F^\ast$ 的作用，采用 Galois 观点最为方便：同构 \eqref{VI:eq:4-12-1} 具有典范性，因而与 $\gal(\dF/\dF_q)$ 经结构迁移所给出的作用相容。若 $\varepsilon(k)$ 是 $\varphi$ 在 $I$ 上的置换符号，则
\begin{align*}
\tau_{\dF_q}(\chi\circ N_{k/\dF_q},\psi)
&=\tr\left(\varphi^{-1},\h_c^N(V^\ast,N^\ast\sF(\chi^{-1})\otimes T^\ast\sF(\psi))\right)\\
&=\varepsilon(k)\tr\left(\varphi^{-1},\h_c^1(\dG_m,\sF(\chi^{-1}\psi))\right)^N\\
&=\varepsilon(k)\tau(\chi,\psi)^N.
\end{align*}
也就是说，
\begin{equation*}\tag{4.12.2}\label{VI:eq:4-12-2}
\tau_{\dF_q}(\chi\circ N_{k/\dF_q},\psi)=\varepsilon(k)\tau(\chi,\psi)^N.
\end{equation*}
若 $k$ 为域，则 $\varphi$ 是 $I$ 的一个循环置换，$\varepsilon(k)=(-1)^{N+1}$。于是得到 Hasse--Davenport 恒等式：
\[
\tau(\chi\circ N_{k/\dF_q},\psi\circ\tr_{k/\dF_q})
=(-1)^{N+1}\tau_{\dF_q}(\chi\circ N_{k/\dF_q},\psi)=\tau(\chi,\psi)^N.
\]

\begin{lemma_}\label{VI:4-13}
令 $\chi\circ\tau=(\chi_i)_{i\in I}$，记号如 \ref{VI:4-9}(i)。以下条件等价：
\begin{enumerate}[\indent (i)]
\item $\chi|\dF_q^\times$ 平凡；
\item 所有 $\chi_i$ 的乘积平凡；
\item $\sF(\chi)=\sK_n((\chi_i))$ 是 $V^\ast/\dG_m$ 上某个唯一的层的逆像；
\item 将 $\sF(\chi)$ 沿由结构态射 $\dF_q\to k$ 给出的 $\dG_m\to V^\ast$ 拉回，所得层平凡。
\end{enumerate}
\end{lemma_}
(i)$\Rightarrow$(ii)：把 $\chi$ 看成 $V_0^\ast/\dG_m(\dF_q)=k^\times/\dF_q^\times$ 的特征，并在 $V_0^\ast/\dG_m$ 上取 $\sF(\chi)$。

(ii)$\Rightarrow$(iii)：类似地，把 $(\chi_i)$ 看成 $(V^\ast/\dG_m)_n$ 的特征。(iii) 中的唯一性来自 $V^\ast$ 是 $V^\ast/\dG_m$ 上纤维连通的纤维丛；(iii)$\Rightarrow$(iv) 则显然。

(iv)$\Rightarrow$(i)、(ii)：该逆像层同时为 $\sF(\chi|\dF_q^\times)$ 和 $\sK_n(\prod\chi_i)$。

\subsection{}\label{VI:4-14}
设 $k$ 是 $\dF_q$ 上维数为 $N+1$ 的平展代数，$\chi$ 是 $k^\times$ 的非平凡特征，且其在 $\dF_q^\times$ 上的限制平凡。\term{Jacobi 和}{Jacobi sum}{somme de Jacobi} $J(\chi)$ 定义为
\begin{equation*}\tag{4.14.1}\label{VI:eq:4-14-1}
J(\chi)=(-1)^{N-1}\sum_{\substack{x\in k^\times/\dF_q^\times\\\tr(x)=0}}\chi^{-1}(x).
\end{equation*}
Gauss 和与 Jacobi 和之间有如下恒等式。

\begin{proposition_}\label{VI:4-15}
若 $\chi$ 如上，$\psi$ 是 $\dF_q$ 的非平凡加法特征，则
\[qJ(\chi)=\tau_{\dF_q}(\chi,\psi).\]
\end{proposition_}
特别地，当 $k=\dF_q^{N+1}$ 时，由 \eqref{VI:eq:4-5-3}，此式写为
\begin{equation*}\tag{4.15.1}\label{VI:eq:4-15-1}
qJ(\chi)=\prod_i\tau(\chi_i,\psi),
\end{equation*}
当 $\chi_0$ 非平凡时，常见的写法是
\begin{align*}
\chi_0(-1)J(\chi)
&=(-1)^{N-1}\sum_{\substack{x_1,\ldots,x_N\in\dF_q^\times\\\sum x_i=1}}\prod_{i=1}^N\chi_i^{-1}(x_i)
=\tau(\chi_0^{-1},\psi)^{-1}\prod_{i=1}^N\tau(\chi_i,\psi),\\
\chi_0^{-1}&=\prod_{i=1}^N\chi_i.
\end{align*}
\translatornote{原书未附加 $\chi_0\ne1$ 这一条件，但最后一个等号需要它。由 (4.15.1)，$qJ(\chi)=\tau(\chi_0,\psi)\prod_{i=1}^N\tau(\chi_i,\psi)$；当 $\chi_0\ne1$ 时，用 $\tau(\chi_0,\psi)\tau(\chi_0^{-1},\psi)=q\chi_0(-1)$ 即得上式。当 $\chi_0=1$ 时，$\tau(1,\psi)=1$，故正确形式是 $J(\chi)=q^{-1}\prod_{i=1}^N\tau(\chi_i,\psi)$。例如取 $q=3$、$N=2$、$\chi_1=\chi_2$ 为二次特征，原书最后一个表达式为 $-3$，而 $J(\chi)=-1$。}
\begin{proof}[命题 \ref{VI:4-15} 的证明]
\begin{align*}
\tau_{\dF_q}(\chi,\psi)
&=(-1)^{N+1}\sum_{x\in k^\times}\chi(x)^{-1}\psi(\tr x)\\
&=(-1)^{N+1}\sum_{x\in k^\times/\dF_q^\times}\chi(x)^{-1}
 \sum_{\lambda\in\dF_q^\times}\psi(\lambda\tr x).
\end{align*}
当 $\tr(x)=0$ 时，内和为 $q-1$；否则为 $-1$。因此
\[
\tau_{\dF_q}(\chi,\psi)=(-1)^{N+1}\left(
 q\sum_{\substack{x\in k^\times/\dF_q^\times\\\tr(x)=0}}\chi(x)^{-1}
 -\sum_{x\in k^\times/\dF_q^\times}\chi(x)^{-1}\right).
\]
右侧第二项是非平凡特征的值之和，故为零，\ref{VI:4-15} 因而成立。
\end{proof}
命题 \ref{VI:4-15} 可提升为如下上同调陈述。

\begin{proposition_}\label{VI:4-16}
$W^\ast/\dG_m$（\ref{VI:4-6}、\ref{VI:4-8}）上以 $\sF(\chi^{-1})$ 为系数的上同调满足：
\begin{enumerate}[\indent (i)]
\item 当 $i\ne N-1$ 时，$\h_c^i=0$，且 $\dim\h_c^{N-1}=1$；
\item 在 $\h_c^{N-1}$ 上，$F^\ast$ 的作用是乘以 $J(\chi)$；
\item $\h_c^{N-1}$ 典范同构于
$\h_c^{N+1}(V^\ast,\sF(\chi^{-1}\cdot\psi\circ\tr_{k/\dF_q}))(1)$。
\end{enumerate}
\end{proposition_}
断言 (ii) 由 (i) 和迹公式推出；而 (i)、(iii) 可由 \ref{VI:4-11} 及下面更具几何性的命题得到。

\begin{proposition_}\label{VI:4-17}
设 $(\chi_i)_{i\in I}$ 是 $\dmu_n(\dF)$ 的一族不全平凡、乘积为 $1$ 的特征。则
$\h_c^{i-1}(W^\ast/\dG_m,\sK_n((\chi_i)_{i\in I}))$
典范同构于
$\h_c^{i+1}(V^\ast,\sK_n((\chi_i)_{i\in I})\otimes\sF(\psi\tr_{k/\dF_q}))(1)$。
\end{proposition_}
\ref{VI:4-15} 的证明第一行成为如下公式（$\pi$ 如 \eqref{VI:eq:4-6-1}）：
\begin{equation*}\tag{4.17.1}\label{VI:eq:4-17-1}
\eR^i\pi_!\left(\sK_n((\chi_i)_{i\in I})\otimes\sF(\psi\tr_{k/\dF_q})\right)
=\sK_n((\chi_i)_{i\in I})\otimes\eR^i\pi_!\sF(\psi\tr_{k/\dF_q}),
\end{equation*}
这是 $V^\ast/\dG_m$ 上的等式。下面计算
$\eR^i\pi_!\sF(\psi\tr_{k/\dF_q})=\eR^i\pi_!T^\ast\sF(\psi)$。
设 $v_0:\widetilde V_0\to V_0$ 是 $V_0$ 在 $\{0\}$ 处的\term{爆破}{blow-up}{éclatement}。在图
\begin{equation*}\tag{4.17.2}\label{VI:eq:4-17-2}
\xymatrix{
 V_0\setminus\{0\}\ar@{^{(}->}[r]\ar[dr]_-\pi&\widetilde V_0\ar[r]^-{v_0}\ar[d]^-{\bar\pi}&V_0\\
 &P_0
}
\end{equation*}
中，$\pi$ 是纤维为去心直线的纤维化，$\widetilde V_0$ 是相应的线丛，$Z_0=v_0^{-1}(0)$ 是其零截面。$V_0\setminus\{0\}$ 上的层 $T^\ast\sF(\psi)$ 可延拓为 $\widetilde V_0$ 上的 $(Tv_0)^\ast\sF(\psi)$。由于 $T_0v_0$ 在 $Z_0$ 上为零，该层在 $Z_0$ 上的限制为常值层 $E_\lambda$。因此固有支撑上同调长正合列为
\begin{equation*}\tag{4.17.3}\label{VI:eq:4-17-3}
\xymatrix{
\ar[r]&\eR^i\pi_!T^\ast\sF(\psi)\ar[r]&\eR^i\bar\pi_!(Tv)^\ast\sF(\psi)\ar[r]&
 {\begin{cases}E_\lambda,&i=0,\\0,&i\ne0\end{cases}}\ar[r]&
}
\end{equation*}
$\bar\pi^{-1}(Q_0)$ 是 $W_0$ 在 $0$ 处的爆破 $\widetilde W_0$，而 $T_0v_0$ 在 $\widetilde W_0$ 上为零。因此在 $\bar\pi^{-1}(Q_0)$ 上有 $(Tv)^\ast\sF(\psi)=E_\lambda$。又因 $\bar\pi$ 为线丛，故
\begin{equation*}\tag{4.17.4}\label{VI:eq:4-17-4}
\eR^i\bar\pi_!(Tv)^\ast\sF(\psi)|Q=
\begin{cases}0,&i\ne2,\\E_\lambda(-1),&i=2.\end{cases}
\end{equation*}
若 $x\in P$ 且 $x\notin Q$，则 $Tv$ 将直线 $D=\bar\pi^{-1}(x)$ 同构地映到 $\dG_a$。因此由 \ref{VI:2-7*}，有
\begin{align}\notag
\h_c^\bullet(D,(Tv)^\ast\sF(\psi))&=\h_c^\bullet(\dG_a,\sF(\psi))=0,\\
\tag{4.17.5}\label{VI:eq:4-17-5}
\eR^i\bar\pi_!(Tv)^\ast\sF(\psi)&=
\begin{cases}0,&i\ne2,\\E_\lambda(-1)_Q,&i=2.\end{cases}
\end{align}
结合 \eqref{VI:eq:4-17-3} 与 \eqref{VI:eq:4-17-5}，最后得到
\begin{equation*}\tag{4.17.6}\label{VI:eq:4-17-6}
\eR^i\pi_!T^\ast\sF(\psi)=
\begin{cases}0,&i\ne1,2,\\E_\lambda,&i=1,\\E_\lambda(-1)_Q,&i=2.\end{cases}
\end{equation*}

\subsection{}\label{VI:4-18}
用 $\pi:V^\ast\to V^\ast/\dG_m$ 的 Leray 谱序列计算 $V^\ast$ 上的层
$\sK_n((\chi_i))\otimes T^\ast\sF(\psi)$ 的固有支撑上同调。应用 \eqref{VI:eq:4-17-1} 与 \eqref{VI:eq:4-17-6}，初始项为
$E_2^{p,1}=\h_c^p(V^\ast/\dG_m,\sK_n((\chi_i)))=0$
（因为 $(\chi_i)\ne1$），以及
$E_2^{p,2}=\h_c^p(W^\ast/\dG_m,\sK_n((\chi_i)))(-1)$。
因此谱序列退化为一个同构，\ref{VI:4-17} 即得证。

\subsection{}\label{VI:4-19}
层 $\sK_n((\chi_i)_{i\in I})$ 在任意域、甚至任意基概形上都有意义。由此可推广 \ref{VI:4-16}(i)。沿用 \ref{VI:4-6} 的记号，设 $M$ 在 $A$ 上有限平展，处处秩为 $N$，且整数 $n$ 在 $A$ 中可逆。记 $a:W^\ast/\dG_m\to\spec(A)$ 为投影。再设
$\chi:(V^\ast/\dG_m)_n\to E_\lambda^\times$
是一个特征，即层的态射，$\sK_n(\chi)$ 为对应的层。在平展拓扑下局部地，可将 $\chi$ 看成 $\dmu_n$ 的一族乘积平凡的特征 $(\chi_i)_{i\in I}$。

\begin{proposition_}\label{VI:4-20}
\begin{enumerate}[(i)]
\item 若 $\chi$ 在每一点均非平凡，则当 $i\ne N-2$ 时，$\eR^ia_!(\sK_n(\chi))=0$，且 $\eR^{N-2}a_!(\sK_n(\chi))$ 是秩 $1$ 的光滑层。
\item $M$ 的一个保持 $\chi$ 的自同构 $\sigma$ 在此 $\eR^{N-2}a_!$ 上的作用，是乘以 $\sigma$ 作为 $I$ 的置换的符号 $\varepsilon(\sigma)$。
\item 若所有 $\chi_i$ 均非平凡，则 $\eR a_!\iso\eR a_\ast$。
\end{enumerate}
\end{proposition_}
\translatornote{原书本命题 (i)、(ii) 将次数写为 $N-1$，应为 $N-2$：这里 $M$ 的秩为 $N$，故 $\dim(W^\ast/\dG_m)=N-2$；在 \ref{VI:4-16} 中秩为 $N+1$，非零次数才是 $N-1$。同篇 \ref{VI:6-2} 也使用 $N-2$。具体反例是 $A=\dF_q$、$M=A^2$，取 $q>2$ 及互为逆的两个非平凡 Kummer 特征。此时 $W^\ast/\dG_m$ 是一个点，非零的紧支撑上同调在零次，而不是一次。将本命题中的 $N-1$ 全部改为 $N-2$ 后，以下证明可直接应用 \ref{VI:4-17}、\ref{VI:4-11}。}
\begin{proof}
$W^\ast/\dG_m$ 是 $Q$ 中一个相对\term{正规交叉除子}{normal crossings divisor}{diviseur à croisements normaux}的补，而 $Q$ 在 $\spec(A)$ 上固有且光滑。层 $\sK_n(\chi)$ 在 $W^\ast/\dG_m$ 上局部常值，在无穷远处温和分歧。因此 $\eR^ia_!$ 和 $\eR^ia_\ast$ 都是光滑层，且其形成与任意基变换相容。一个标准论证将问题归约到 $A$ 是有限域的情形，(i) 随即由 \ref{VI:4-17} 和 \ref{VI:4-11}(i) 得到。为证明 (ii)，只需使用 $\sigma$ 的作用与同构 \ref{VI:4-17} 的相容性，再应用 \ref{VI:4-12}。对于 (iii)，注意若所有 $\chi_i$ 都非平凡，则 $W^\ast/\dG_m\subset Q$ 上的 $\sK_n(\chi)$ 沿无穷远处的每个除子都分歧，再应用 \ref{VI:1-19-1} 即可。
\end{proof}

\section{Hecke 特征}\label{VI:5}
\origpage{214}{208}

\subsection{}\label{VI:5-1}
设 $F$ 是数域，即 $\dQ$ 的有限次扩张，$k$ 是特征为 $0$ 的域。对于同态 $\alpha:F^\times\to k^\times$，可用以下几种等价方式表达它是\term{代数同态}{algebraic homomorphism}{homomorphisme algébrique}。
\begin{enumerate}[(i)]
\item 设 $\{e_i\}$ 是 $F$ 在 $\dQ$ 上的一组基。若 $\alpha$ 可由公式
$\alpha(\sum x^ie_i)=A(x^i)$，$A\in k(X^i)$ 给出，则称其为代数同态。
\end{enumerate}
这意味着 $\alpha$ 在 $F^\times$ 上与一个在 $k$ 上有定义的有理映射
$\res_{F/\dQ}(\dG_m)\dashrightarrow\dG_m$ 一致。
由于 $F^\times$ Zariski 稠密，且有理群同态处处有定义，这样的映射是群概形的同态：
\begin{enumerate}[(i)]\setcounter{enumi}{1}
\item $\alpha$ 由一个 $k$-群概形同态
\[\res_{F/\dQ}(\dG_m)\otimes_\dQ k\longrightarrow\dG_m\]
诱导。
\end{enumerate}
若 $k$ 为数域，则\term{标量限制}{restriction of scalars}{restriction des scalaires} $R$ 的伴随性质表明，这又等价于：
\begin{enumerate}[(i)]\setcounter{enumi}{2}
\item （$k$ 为数域。）$\alpha$ 由 $\dQ$-群概形同态
$\res_{F/\dQ}(\dG_m)\to\res_{k/\dQ}(\dG_m)$ 诱导。
\end{enumerate}
设 $\bar k$ 为 $k$ 的代数闭包，$I=\hom(F,\bar k)$。在 $\bar k$ 上，$\res_{F/\dQ}(\dG_m)$ 的特征群是 $\dZ^I$，以 $F$ 到 $\bar k$ 的各嵌入为基。在 $k$ 上有定义的特征恰好是 $\gal(\bar k/k)$ 不变的特征；既可以将其描述为把 $F^\times$ 映入 $k^\times\subset\bar k$ 的特征，也可以利用 $\gal(\bar k/k)$ 在 $I$ 中的轨道来描述。这些轨道又对应于 $F\otimes_\dQ k$ 的各因子。
\begin{enumerate}[(i)]\setcounter{enumi}{3}
\item $\alpha$ 具有形式 $\alpha=\prod_{\omega\in I}\omega^{n_\omega}$。允许的指数族 $(n_\omega)$ 满足：若 $\omega,\omega'$ 位于同一 $\gal(\bar k/k)$-轨道，则 $n_\omega=n_{\omega'}$。等价地，对每个 $x\in F^\times$，有 $\prod_\omega\omega(x)^{n_\omega}\in k$。
\item 写 $F\otimes_\dQ k=\prod_{j\in J}F_j$，其中各 $F_j$ 为域，则 $\alpha$ 可写为 $\alpha=\prod_jN_{F_j/k}^{m_j}$。
\end{enumerate}
特别地，若 $k$ 包含 $F$ 的一个正规闭包，则这些表达式成为 $\alpha=\prod\omega^{n_\omega}$，其中 $\omega$ 遍历 $F$ 到 $k$ 的 $[F:\dQ]$ 个嵌入。

\subsection{}\label{VI:5-2}
设 $k$ 为数域，$\alpha:F^\times\to k^\times$ 为代数同态。那么恰有一个同态，仍记作 $\alpha$，从 $F$ 的\term{分式理想}{fractional ideal}{idéal fractionnaire}群 $I(F)$ 到 $k$ 的分式理想群，满足 $\alpha((x))=(\alpha(x))$。唯一性来自每个理想都有一个幂为主理想，且 $I(k)$ 无挠。为证存在性，例如可利用 \ref{VI:5-1}(v)：令
$\alpha(\mathfrak a)=\prod_jN_{F_j/k}^{m_j}(\mathfrak a\mathcal O_{F_j})$ 即可。

\subsection{}\label{VI:5-3}
回顾\term{代数 Hecke 特征}{algebraic Hecke character}{caractère de Hecke algébrique}的定义。Weil 称其为 $A_0$ 型 Hecke 特征，亦称 Grössencharakter。设 $F$ 为数域，$\fm$ 为 $F$ 的一个整理想，即 $F$ 的整数环的理想；$I_\fm$ 是 $F$ 中与 $\fm$ 互素的分式理想群，$k$ 为特征 $0$ 的域。若同态 $\chi:I_\fm\to k^\times$ 满足：存在代数同态 $\chi_\text{alg}:F^\times\to k^\times$，使得
\begin{equation*}\tag{$*$}\label{VI:eq:5-3-1}
\begin{gathered}
 x\in F^\times\text{ 与 }\fm\text{ 互素、全正，且 }x\equiv1\pmod\fm\\
 \Longrightarrow\quad\chi((x))=\chi_\text{alg}(x),
\end{gathered}
\end{equation*}
则称 $\chi$ 是一个代数 Hecke 特征，其\term{导子}{conductor}{conducteur} $\le\fm$。

由 \eqref{VI:eq:5-3-1} 中这些 $x$ 的稠密性，$\chi_\text{alg}$ 完全由 $\chi$ 决定，称为 $\chi$ 的\term{代数部分}{algebraic part}{partie algébrique}。若对于全正且 $x\equiv1\pmod{\fm'}$ 的 $x$，有 $\chi((x))=\chi_\text{alg}(x)$，则 $\chi$ 可延拓为导子 $\le\inf(\fm,\fm')$ 的 Hecke 特征 $I_{\fm+\fm'}\to k^\times$。在公共定义域上相同的 Hecke 特征视为同一特征；使 $\chi$ 的导子 $\le\fm'$ 的最小 $\fm'$，称为 $\chi$ 的导子。

\subsection{备注}\label{VI:5-4}
若代数 Hecke 特征 $\chi$ 的值落在 $k$ 的子域 $k'$ 中，则它本来就是一个 $k'$ 值的 Hecke 特征。只需说明 $\chi_\text{alg}:\res_{F/\dQ}(\dG_m)\to\dG_m$ 在 $k'$ 上已有定义，而这来自满足 $x\equiv1\pmod\fm$ 的全正元在 $\res_{F/\dQ}(\dG_m)$ 中的 Zariski 稠密性。总可以将 $k'$ 取为 $k$ 中在 $\dQ$ 上次数有限的子域。

\subsection{}\label{VI:5-5}
若 $\varepsilon$ 为全正单位，且 $\varepsilon\equiv1\pmod\fm$，则
$\chi_\text{alg}(\varepsilon)=\chi((\varepsilon))=1$。
因此，$\chi_\text{alg}$ 通过 $\res_{F/\dQ}(\dG_m)$ 的商 $T_\fm$ 分解；这里的商是对全正且 $\equiv1\pmod\fm$ 的单位群 $E_\fm\subset F^\times$ 的 Zariski 闭包取商。

由 Serre \cite[II.3]{se68}，若 $k$ 是 $\dQ$ 的代数闭包，且 $\fm$ 足够大，则能够作为 $\chi_\text{alg}$ 出现的 $\res_{F/\dQ}(\dG_m)$ 的特征 $\prod\omega^{n_\omega}$，恰可刻画如下：应存在一个整数 $N$，称为 $\chi$ 或 $\chi_\text{alg}$ 的\term{权}{weight}{poids}，使得对于 $\gal(k/\dQ)$ 中每个与复共轭共轭的元素 $\sigma$，都有 $n_\omega+n_{\sigma\omega}=N$。若 $F_1$ 是 $F$ 中最大的、可写成全实域 $F_1'$ 的全虚二次扩张的子域，则这一条件等价于 $\chi_\text{alg}=\chi_1\circ N_{F/F_1}$，其中 $\chi_1|F_1'=(N_{F_1'/\dQ})^N$。

若 $\chi$ 的权为 $N$，则对每个与导子互素的 $F$ 的理想 $\fa$，$\chi(\fa)$ 是代数数，其所有共轭的绝对值均为 $N(\fa)^{N/2}$ \cite[II.3，命题 2]{se68}。

\subsection{}\label{VI:5-6}
固定数域 $k$ 后，$\chi_\text{alg}$ 还须满足一些额外条件。事实上，对 $F$ 的每个与导子互素的理想 $\fa$，都有
\begin{equation*}\tag{5.6.1}\label{VI:eq:5-6-1}
\chi_\text{alg}(\fa)=(\chi(\fa)),
\end{equation*}
故 $\chi_\text{alg}(\fa)$ 必须是主理想。由于 $k$ 的分式理想群无挠，只需对某个合适的非零整数 $n$，证明 \eqref{VI:eq:5-6-1} 的 $n$ 次幂。这样便可用 $\fa^n$ 代替 $\fa$，从而假设 $\fa=(x)$，其中 $x$ 全正且 $x\equiv1\pmod\fm$。此时只需使用定义即可。

结合 \ref{VI:5-5}，可知 $\chi_\text{alg}$ 决定了 $\chi(\fa)$ 在 $k$ 的所有位处的绝对值。

\subsection{}\label{VI:5-7}
Serre 群 $S_\fm$ 可刻画为如下\term{乘法型群}{group of multiplicative type}{groupe de type multiplicatif}：其在任意域上的特征群，都是导子 $\le\fm$ 的代数 Hecke 特征群（参见 \cite[II 2.1、2.2]{se68}）。它与 $\ell$-进表示的关系见 \cite[II 2.3]{se68}。

\begin{theorem_}[{\cite{se68}}]\label{VI:5-8}
设 $\chi$ 是 $F$ 的一个取值于 $E_\lambda$ 的代数 Hecke 特征，其中 $E_\lambda$ 是 $\dQ_\ell$ 的有限扩张。那么存在唯一的连续同态
$\chi_\lambda:\gal(\bar F/F)^\textnormal{ab}\to E_\lambda^\times$，使得
\begin{enumerate}[\indent (i)]
\item $\chi_\lambda$ 在 $\chi$ 的导子 $\ff$ 和 $\ell$ 之外不分歧；
\item 若 $\fp$ 是 $F$ 中与 $\ff$ 和 $\ell$ 互素的素理想，且 $F_\fp\in\gal(\bar F/F)^\textnormal{ab}$ 为 $\fp$ 处的几何 Frobenius，则
\[\chi_\lambda(F_\fp)=\chi(\fp).\]
\end{enumerate}
\end{theorem_}

\subsection{}\label{VI:5-9}
在本节余下部分，我们给出一个判据，用来判断 $\lambda$-进表示是否以这种方式来自 Hecke 特征；下一节则将它用于 Jacobi 和。由此将重新得到 Weil 的结论 \cite{we52,we74-2}，并略作推广；其中 Weil 证明的后半部分将由一个 $\ell$-进上同调论证替代。

\begin{theorem_}\label{VI:5-10}
设 $F,k$ 为数域，$\lambda$ 是 $k$ 的剩余特征为 $\ell$ 的一个位，且
$\chi_\lambda:\gal(\bar F/F)^\textnormal{ab}\to k_\lambda^\times$
是连续同态，在 $\ell$ 和有限位集 $S$ 之外不分歧。假设存在密度为 $0$ 的位集 $T$，包含 $S$ 以及 $\ell$ 上方的各位，并存在代数同态 $\chi_0:F^\times\to k^\times$，使得
\begin{enumerate}[(i)]
\item 对 $k$ 的代数闭包 $\bar k$，$\chi_0:F^\times\to k^\times\hookrightarrow\bar k^\times$ 属于 \ref{VI:5-5} 所考虑的类型，权为 $N$；
\item 对每个不在 $T$ 中的素理想 $\fp$，$\chi_\lambda(F_\fp)$ 属于 $k^\times$，且 $(\chi_\lambda(F_\fp))=\chi_0(\fp)$，并且其所有复共轭的绝对值均为 $(N\fp)^{N/2}$。
\end{enumerate}
那么 $\chi_\lambda$ 由 $F$ 的一个 $k$ 值代数 Hecke 特征按 \ref{VI:5-8} 定义，其代数部分为 $\chi_0$。
\end{theorem_}
设 $\chi'$ 是一个取值于 $k$ 的有限 Galois 扩张 $k'$ 的代数 Hecke 特征，其代数部分为 $\chi_0$（\ref{VI:5-5}）。适当扩大 $T$ 后，可假设 $\chi'$ 的导子支撑在 $T$ 中。设 $\lambda'$ 是 $k'$ 中位于 $\lambda$ 上方的一个位，$\chi_{\lambda'}$ 为复合
$\gal(\bar F/F)^\text{ab}\to k_\lambda^\times\to k'_{\lambda'}{}^\times$，而
$\chi'_{\lambda'}:\gal(\bar F/F)^\text{ab}\to k'_{\lambda'}{}^\times$
为 $\chi'$ 按 \ref{VI:5-8} 定义的特征。令
$\varepsilon=\chi_{\lambda'}\cdot(\chi'_{\lambda'})^{-1}$。

假设 (ii) 以及 \ref{VI:5-5}、\ref{VI:5-6}、\ref{VI:5-8} 保证，对于 $\fp\notin T$，有 $\varepsilon(F_\fp)\in k'^\times$，且这个数在 $k'$ 的每个完备化中的绝对值均为 $1$。因此 $\varepsilon(F_\fp)$ 属于 $k'$ 的有限单位根群 $\dmu$。由 Chebotarev 密度定理及 $T$ 的密度假设，$F_\fp$（$\fp\notin T$）在 $\gal(\bar F/F)^\text{ab}$ 中稠密，故对每个 $\sigma\in\gal(\bar F/F)^\text{ab}$，均有 $\varepsilon(\sigma)\in\dmu$。于是 $\varepsilon$ 是有限阶特征 $\gal(\bar F/F)^\text{ab}\to\dmu$。根据类域论，它对应于 $F$ 的 idèle 类群的一个有限阶特征，即代数部分平凡的代数 Hecke 特征。用它修正 $\chi'$ 后，可假设 $\chi_{\lambda'}=\chi'_{\lambda'}$，只需再证明原先取值于 $k'$ 的 $\chi'$ 实际取值于 $k$。

若 $\sigma\in\gal(k'/k)$，则 $\chi'$ 和 $(\chi')^\sigma$ 在每个 $\fp\notin T$ 上一致，因为
$\chi'(\fp)=\chi'_{\lambda'}(F_\fp)=\chi_{\lambda'}(F_\fp)\in k^\times$。
于是对应的 $\lambda'$-进特征在 $\gal(\bar F/F)^\text{ab}$ 的一个稠密子集上一致，从而相等，故 $\chi'=(\chi')^\sigma$。因此 $\chi'$ 取值于 $k$。

\section{由 Jacobi 和定义的 Hecke 特征}\label{VI:6}
\origpage{219}{213}
本节的思路如下。考虑一个 Gauss 和 $\tau_{\dF_q}(\chi,\psi)$。这个和以及相应的层 $\sF(\chi^{-1}\psi)$，通过以下两点与特征 $p$ 及基域 $\dF_q$ 相联系：
\begin{enumerate}[\indent a)]
\item 构造中出现了加法特征 $\psi$ 及层 $\sF(\psi)$；
\item $\sF(\chi)$ 由 Lang 正合列
\[\xymatrix{0\ar[r]&k^\times\ar[r]&V^\ast\ar[r]&V^\ast\ar[r]&0}\]
定义。
\end{enumerate}
若想把相应的上同调构造提升到数域 $F$，并找到 $\gal(\bar F/F)$ 的 $\ell$-进表示，使其 Frobenius 特征值为 Gauss 和，就必须克服这两点困难。

若 $\chi$ 在 $\dF_q^\times$ 上平凡，则 $\tau_{\dF_q}(\chi,\psi)$ 与 $\psi$ 无关，且等于某个 Jacobi 和的 $q$ 倍。这个 Jacobi 和对应于 $W^\ast/\dG_m$ 上以由 $\chi$ 导出的层为系数的上同调；$\psi$ 以及困难 a) 都消失了。为解决 b)，只需借助在任意特征下都有意义的 Kummer 正合列，描述所用的层。由 \ref{VI:4-9}，这是可行的。

找到这些 $\ell$-进表示后，再利用 \ref{VI:5-10} 和 Stickelberger 定理，证明它们来自代数 Hecke 特征。

\subsection{}\label{VI:6-1}
若 $E$ 是域，$\bar E$ 为其代数闭包，而 $\lambda$ 是 $\mu_d(\bar E)$ 的一个阶数整除 $n$ 的特征，则以 ${}_n\lambda$ 表示 $\mu_n(\bar E)$ 的如下特征：$\lambda$ 和 ${}_n\lambda$ 在
$\widehat\dZ(1)_{\bar E}=\varprojlim_m\mu_m(\bar E)$
上诱导相同的特征。若 $\lambda$ 本身是 $\widehat\dZ(1)_{\bar E}$ 的特征，也采用相同的记号。

设 $F,k$ 为两个数域，$\bar F$ 为 $F$ 的代数闭包，$I$ 是带有 $\gal(\bar F/F)$ 作用的有限集，$\lambda=(\lambda_i)_{i\in I}$ 是一族特征
$\lambda_i:\widehat\dZ(1)_{\bar F}\to k^\times$。
假设每个 $\lambda_i\ne1$，它们的乘积平凡，且族 $\lambda$ 在 $F$ 上有定义，即对于所有 $\sigma\in\gal(\bar F/F)$，有
$\lambda_{\sigma(i)}=\lambda_i\circ\sigma^{-1}$。

Galois 集合 $I$ 是某个可分 $F$-代数 $E$ 到 $\bar F$ 的同态集。若 $C$ 为 $\gal(\bar F/F)$ 在 $I$ 中的轨道集，则有域的直积分解
$E=\prod_{\alpha\in C}E_\alpha$，其中 $\alpha$ 等同于 $E_\alpha$ 到 $\bar F$ 的嵌入集。若 $i\in\alpha$ 时各 $\lambda_i$ 的阶为 $d_\alpha$，则 $E_\alpha$ 包含全部 $d_\alpha$ 次单位根，并存在
$\lambda_\alpha:\mu_{d_\alpha}(E_\alpha)\hookrightarrow k^\times$，使得当 $i\in\alpha$ 对应于嵌入 $\sigma_i:E_\alpha\hookrightarrow\bar F$ 时，有
${}_{d_\alpha}\lambda_i=\lambda_\alpha\circ\sigma_i^{-1}$。

\subsection{}\label{VI:6-2}
设 $n>0$ 是所有 $d_\alpha$ 的公倍数，$\Sigma$ 是 $F$ 中 $E/F$ 分歧的位与整除 $n$ 的位所成的集合；$\mu$ 为 $k$ 的一个剩余特征为 $\ell$ 的位；$A$ 为 $F$ 中在 $\Sigma$ 和 $\ell$ 之外整的元素所成的环。$A$ 在 $E$ 中的整闭包 $M$ 在 $A$ 上有限平展，秩为 $N=|I|$。沿用 \ref{VI:4-6}、\ref{VI:4-8} 的记号，各 ${}_n\lambda_i$ 定义了
\[{}_n\lambda:(V^\ast/\dG_m)_n\longrightarrow k^\times\subset k_\mu^\times.\]
设 $a:W^\ast/\dG_m\to\spec(A)$ 为投影。由 \ref{VI:4-20}，层
$\eR^{N-2}a_!\sK_n({}_n\lambda)$ 光滑且秩为一，故定义了 $\ell$-进表示
\[j_\mu[\lambda]:\gal(\bar F/F)^\text{ab}\longrightarrow k_\mu^\times,\]
也简记为 $j_\mu$，它在 $\Sigma$ 和 $\ell$ 之外不分歧。

\subsection{}\label{VI:6-3}
下面计算 $j_\mu(F_\fp)$。由 \ref{VI:4-20}，只需在模 $\fp$ 约化后计算。令 $f=A/\fp$，$e=M/\fp M$，并设 $\bar f$ 是由 $\bar F$ 中位于 $\fp$ 上方的一个位所确定的 $f$ 的代数闭包。仍有 $I=\hom_f(e,\bar f)$。约化后，${}_n\lambda$ 仍可按 \ref{VI:6-1} 描述：
\begin{enumerate}[a)]
\item 设 $D$ 为 $\gal(\bar f/f)$ 在 $I$ 中的轨道集。将 $e$ 分解为域的直积 $e=\prod_{\beta\in D}e_\beta$，则 $\sigma_i$（$i\in\beta$）对应于 $e_\beta$ 到 $\bar f$ 的各嵌入。对 $\beta\in D$，记 $\alpha(\beta)$ 为 $C$ 中包含 $\beta$ 的元素，并令 $d_\beta=d_{\alpha(\beta)}$。令 $\lambda_\beta$ 为复合
\[\xymatrix{
\mu_{d_\beta}(e_\beta)&\mu_{d_\beta}(E_{\alpha(\beta)})\ar[l]_-\sim\ar[r]^-{\lambda_{\alpha(\beta)}}&k^\times.
}\]
\item 若 $i\in I$ 对应于嵌入 $\sigma_i:e_\beta\hookrightarrow\bar f$，则 ${}_{d_\beta}\lambda_i$ 的约化
${}_{d_\beta}\bar\lambda_i:\mu_{d_\beta}(\bar f)\to k^\times$
等于 $\lambda_\beta\circ\sigma_i^{-1}$。
\end{enumerate}
设 $q_\beta=|e_\beta|$，$\chi_\beta={}_{q_\beta-1}\lambda_\beta$，$\chi:e^\times\to k^\times$ 是以各 $\chi_\beta$ 为分量的特征。在由 $e$ 定义的 $A/\fp$ 上的仿射空间（\ref{VI:4-6}）上，更确切地说，在环面 $e_\beta^\times$ 上，有同构
\begin{align*}
\sK_n(({}_n\bar\lambda_i)_{i\in\beta})
&=\sK_{d_\beta}((\lambda_\beta\circ\sigma_i^{-1})_{i\in\beta})\\
&=\sK_{q_\beta-1}((\chi_\beta\circ\sigma_i^{-1})_{i\in\beta})
=\sF(\chi_\beta)
\end{align*}
（\ref{VI:4-9}）。因此在 $e^\times$ 上有
$\sK_n(({}_n\bar\lambda_i)_{i\in I})=\sF(\chi)$，且该同构下降到 $V^\ast/\dG_m$ 的模 $\fp$ 约化上。于是得到：

\begin{proposition_}\label{VI:6-4}
沿用 \ref{VI:6-2}、\ref{VI:6-3} 的记号，$j_\mu(F_\fp)$ 就是 Jacobi 和 $J(\chi)$。
\end{proposition_}
将联系 $(\lambda_i)_{i\in I}$ 与 $(\chi_\beta)_{\beta\in D}$ 的公式重写如下。对于 $i\in\beta\subset\alpha$，有
\begin{align*}
{}_{d_\alpha}\lambda_i&=\lambda_\alpha\circ\sigma_i^{-1},
&&\lambda_\alpha:\mu_{d_\alpha}(E_\alpha)\to k^\times;\\
\lambda_\beta&=\lambda_\alpha\text{ 的约化},
&&\lambda_\beta:\mu_{d_\alpha}(e_\beta)\to k^\times;\\
\chi_\beta&={}_{q_\beta-1}\lambda_\beta:
 x\mapsto\lambda_\beta(x^{(q_\beta-1)/d_\alpha}),
&&\chi_\beta:e_\beta^\times\to k^\times.
\end{align*}

\begin{theorem_}\label{VI:6-5}
表示 $j_\mu$ 由 $F$ 的一个 $k^\times$ 值代数 Hecke 特征定义。
\end{theorem_}
我们将利用 \ref{VI:5-10} 证明这一点。在使用 Stickelberger 定理这一环节，我们与 Weil 的证明 \cite{we52,we74-2} 会合；这里只给出若干提示。

\subsection{}\label{VI:6-6}
给定 $j:\gal(\bar F/F)\to k_\lambda^\times$ 及有限扩张 $F'/F$，用 \ref{VI:5-10} 可验证：$j$ 由 Hecke 特征定义，当且仅当其限制 $j|\gal(\bar F/F')$ 如此。由此，在 \ref{VI:6-5} 中可归约到 Galois 群在 $I$ 上作用平凡且 $F$ 包含全部 $n$ 次单位根的情形。进一步可归约到 $F=\dQ(\sqrt[n]{1})$，并取 $k=F$。这时每个 ${}_n\lambda_i$ 由某个 $a_i\in\dZ/n$ 定义，即
$x\mapsto x^{a_i}:\mu_n\to k^\times$，且 $\sum a_i=0$、$a_i\ne0$。

\subsection{}\label{VI:6-7}
在数域中，次数为 $1$ 的位总构成密度为 $1$ 的集合。因此应用 \ref{VI:5-10} 时，可以略去其余各位，只对素域使用 Stickelberger 定理。
\translatornote{原书此处“密度 $0$”应为“密度 $1$”，这也与紧接着的“可以略去其余各位”一致。证明如下：在次数为 $d$ 的数域中，次数至少为二的素理想的 Dirichlet 级数在 $s=1$ 处被 $d\sum_p p^{-2}$ 控制，因而 Dirichlet 密度为零；所有有限位的密度为一，所以一次位的密度为一。最简单的反例是数域 $\dQ$，其每个有限位的次数都为一。}
设 $p$ 为素数，$0<a<p-1$，$\widetilde x:\dF_p^\times\to\dQ_p^\times$ 为“乘法提升”特征，并令
$g=-\sum_{x\in\dF_p^\times}\widetilde x^{-a}\zeta^x\in\dQ_p(\zeta)$，其中 $\zeta$ 为一个非平凡的 $p$ 次单位根。仍用 $x$ 表示 $[0,p-1]$ 中提升 $x$ 的整数，并令 $\pi=\zeta-1$，用二项式公式展开 $\zeta^x=(1+\pi)^x$。得到
$g=-\sum_{i=0}^{p-1}A_i\pi^i$，其中 $A_i\in\dZ_p$，其模 $p$ 约化为 $\sum_{x\in\dF_p^\times}x^{-a}\binom{x}{i}$。若 $b$ 不被 $p-1$ 整除，则
$\sum_{x\in\dF_p^\times}x^b=0$。当 $i<a$ 时，$\binom{x}{i}$ 为次数 $i<a$ 的 $x$ 的多项式，故 $A_i\equiv0\pmod p$。同理，$A_a\equiv(p-1)/a!\pmod p$，于是 $g\sim\pi^a/a!$，即 $g$ 的赋值满足
$v(g)/v(p)=a/(p-1)$。

\subsection{}\label{VI:6-8}
记 $\sigma_i$（$i\in(\dZ/n)^\times$）为 $\gal(\dQ(\sqrt[n]{1})/\dQ)$ 中满足 $\sigma_i(\zeta)=\zeta^i$（$\zeta^n=1$）的元素，并用加法记号表示代数同态群
$\dQ(\sqrt[n]{1})^\times\to\dQ(\sqrt[n]{1})^\times$。
设 $\lambda$ 如 \ref{VI:6-6}，由一族非零且和为零的模 $n$ 整数 $(a_i)_{i\in I}$ 定义。由 \ref{VI:5-10}、\eqref{VI:eq:4-15-1} 和 \ref{VI:6-7}，$j_\mu[\lambda]$ 由一个代数 Hecke 特征 $j[a]$ 定义。若 $N$ 表示范数，则其代数部分由下式给出：
\begin{equation*}\tag{6.8.1}\label{VI:eq:6-8-1}
(Nj[a])_\text{alg}
=\sum_{i\in(\dZ/n)^\times}\sum_{j\in I}
 \left\{\frac{ia_j}{n}\right\}\sigma_i^{-1},
\end{equation*}
其中 $\{\cdot\}$ 表示小数部分。

\subsection{}\label{VI:6-9}
最有意思的情形是 $F\subset\dQ(\sqrt[n]{1})$，且族 $({}_n\lambda_i)$ 由以下步骤构造：
\begin{enumerate}[a)]
\item 在 $\dZ/n$ 中取 $\gal(\dQ(\sqrt[n]{1})/F)=H\subset(\dZ/n)^\times$ 的有限族轨道 $(A_j)_{j\in J}$。假设 $A_j\ne\{0\}$，且 $\sum_{j\in J}\sum_{a\in A_j}a=0$；
\item 取 $I=\bigsqcup_{j\in J}A_j$，Galois 群在每个 $A_j$ 上作用；
\item 取 $\bar F\supset\dQ(\sqrt[n]{1})$，$k=\dQ(\sqrt[n]{1})$；若 $i\in I$ 对应于 $a\in A_j\subset\dZ/n$，则令 ${}_n\lambda_i(x)=x^a$。
\end{enumerate}
用 \ref{VI:4-20} 可验证，一般的 $j_\mu$（相对于 $F_1$）可由上述某个 $j_\mu$ 得到：取 $F\subset\dQ(\sqrt[n]{1})$，$F_1$ 为 $F$ 的扩张，将表示限制到 $\gal(\bar F/F_1)$，再乘以一个二阶特征。

对 $a\in\gal(k/\dQ)=(\dZ/n)^\times$，$a$ 将族 $(A_j)$ 变成 $(aA_j)$。若 $a\in H$，则 $(A_j)$ 被保持，因此各 $j_\mu(F_\fp)$ 均为 $H$ 不变元。

\begin{proposition_}\label{VI:6-10}
在 \ref{VI:6-9} 的假设下，$j_\mu$ 由 $F$ 的一个取值于 $F$ 的代数 Hecke 特征定义。
\end{proposition_}
重新取 $F=\dQ(\sqrt[n]{1})$，并考虑由满足 $a_i\in\dZ/n$、$a_i\ne0$、$\sum a_i=0$ 的族 $(a_i)_{i\in I}$ 定义的 Hecke 特征 $j[a]$。以下定理由 Vishik 独立得到。

\begin{theorem_}\label{VI:6-11}
$\dQ(\sqrt[n]{1})$ 的任意代数 Hecke 特征，都有一个幂属于由各 $j[a]$ 和范数生成的群。
\end{theorem_}
由于我们在“模去挠元”的意义下工作，只有这些 Hecke 特征的代数部分起作用。对 $a$ 由 $\dZ/n$ 中同一个元素重复 $n$ 次而成的 $j[a]$ 应用 \eqref{VI:eq:6-8-1}，即可归约到以下引理。

\begin{lemma_}\label{VI:6-12}
将群代数 $\dQ[(\dZ/n)^\times]$ 的元素记作 $\sum x_j\sigma_j$，令 $X$ 为满足 $x_j+x_{-j}$ 是与 $j$ 无关的常数的那些元素所成的子群。那么 $X$ 在 $\dQ$ 上由 $N=\sum\sigma_j$ 和
\[g_a=\sum_{i\in(\dZ/n)^\times}\left\{\frac{ia}{n}\right\}\sigma_i^{-1}
\qquad(a\in\dZ/n,\ a\ne0)\]
生成。
\end{lemma_}
$X$ 和 $Y=\langle N,(g_a)\rangle$ 都是群代数的理想，因此只需验证它们被同样的复特征 $\chi$ 消去。这等价于说，奇特征，即 $\chi(j)=-\chi(-j)$，不可能将所有 $g_a$ 同时消去。若 $\chi$ 为本原特征，则
\[\chi(g_1)=\sum_{i\in(\dZ/n)^\times}\frac{i}{n}\chi^{-1}(i)
=-L(\chi^{-1},0)\ne0,\]
其中 $i$ 取 $1,\ldots,n-1$ 中的代表元。一般地，若 $\chi$ 来自 $(\dZ/d)^\times$ 的本原特征，则 $\chi(g_{n/d})$ 归结为 $(\dZ/d)^\times$ 上的同类和，故引理成立。

\section{广义 Kloosterman 和}\label{VI:7}
\origpage{225}{219}

\subsection{}\label{VI:7-1}
设 $\dF_q$ 为含 $q$ 个元素的有限域，$\psi$ 为非平凡加法特征。本节从上同调角度研究\term{广义 Kloosterman 和}{generalized Kloosterman sum}{somme de Kloosterman généralisée}
\begin{equation*}\tag{7.1.1}\label{VI;eq:7-1-1}
\kloos_{n,a}=\sum_{x_1\cdots x_n=a}\psi(x_1+\cdots+x_n)
\qquad(n\ge1).
\end{equation*}
当 $a=0$ 时，此和可用初等方法求出（参见 \ref{VI:7-7}）：
\begin{equation*}\tag{7.1.2}\label{VI:eq:7-1-2}
\kloos_{n,0}=(-1)^{n-1}.
\end{equation*}
有意思的是 $a\ne0$ 的情形，此时各 $x_i\in\dF_q^\times$。我们将证明估计
\begin{equation*}\tag{7.1.3}\label{VI:eq:7-1-3}
|\kloos_{n,a}|\le nq^{(n-1)/2}.
\end{equation*}
主要工具是反映下列恒等式的上同调陈述：
\begin{align*}\tag{7.1.4}\label{VI:eq:7-1-4}
\kloos_{n,a}
&=\sum_{x_1\in\dF_q^\times}\psi(x_1)
 \sum_{x_2\cdots x_n=a/x_1}\psi(x_2+\cdots+x_n)\\\notag
&=\sum_{x\in\dF_q^\times}\psi(x)\kloos_{n-1,a/x}
\qquad(a\ne0,\ n\ge2);\\
\tag{7.1.5}\label{VI:eq:7-1-5}
\sum_a\kloos_{n,a}
&=\sum_a\sum_{x_1\cdots x_n=a}\psi\left(\sum x_i\right)
=\sum_{x_1,\ldots,x_n}\psi\left(\sum x_i\right)=0.
\end{align*}
若 $\chi$ 为 $\dF_q^\times$ 的非平凡特征，则
\begin{equation*}\tag{7.1.6}\label{VI:eq:7-1-6}
\sum_a\chi(a)\kloos_{n,a}
=\sum_{x_1,\ldots,x_n}\chi\left(\prod x_i\right)\psi\left(\sum x_i\right)
=(-\tau(\chi^{-1},\psi))^n,
\end{equation*}
右侧是 Gauss 和的幂。

\subsection{}\label{VI:7-2}
设 $k$ 为 $\dF_q$ 上次数为 $n$ 的平展代数，$\varepsilon(k)$ 如 \eqref{VI:eq:4-5-1}。令
\[\kloos_{k,a}=\sum_{N_{k/\dF_q}(x)=a}\psi(\tr(x)).\]
仍有
\begin{align*}\tag{7.2.1}\label{VI:eq:7-2-1}
\kloos_{k,0}&=(-1)^{n-1}\varepsilon(k);\\
\tag{7.2.2}\label{VI:eq:7-2-2}
\sum_a\kloos_{k,a}&=0;\\
\tag{7.2.3}\label{VI:eq:7-2-3}
\sum_a\chi(a)\kloos_{k,a}
&=(-1)^n\tau_{\dF_q}(\chi^{-1}\circ N_{k/\dF_q},\psi).
\end{align*}
在 $\dF_q^\times$ 上应用 Fourier 反演，可将 $\kloos_{k,a}$ 写成 Gauss 和的表达式：
\begin{align*}
\kloos_{k,a}&=\frac1{q-1}\sum_\chi\chi(a)
 \sum_{b\in\dF_q^\times}\chi^{-1}(b)\kloos_{k,b};\\
\tag{7.2.4}\label{VI:eq:7-2-4}
(-1)^n\kloos_{k,a}
&=\frac1{q-1}\left(\varepsilon(k)+\sum_{\chi\ne1}\chi(a)
 \tau_{\dF_q}(\chi\circ N_{k/\dF_q},\psi)\right).
\end{align*}
现在借助 Hasse--Davenport 恒等式，可以将 $\kloos_{k,a}$ 与
$\kloos_{n,a}=\kloos_{\dF_q^n,a}$ 比较：
\begin{equation*}\tag{7.2.5}\label{VI:eq:7-2-5}
\kloos_{k,a}=\varepsilon(k)\kloos_{n,a}.
\end{equation*}
我不知道 \eqref{VI:eq:7-2-4} 的上同调解释，但下面将给出 \eqref{VI:eq:7-2-5} 的上同调解释。先回到和式 $\kloos_{n,a}$。

\subsection{}\label{VI:7-3}
设 $\dF$ 为 $\dF_q$ 的代数闭包。对 $a\in\dF$，用 $V_a^{n-1}$，或简记 $V_a$，表示 $\dA^n$ 中方程 $x_1\cdots x_n=a$ 定义的超曲面。照例，我们将 $\psi$ 看成 $\ell$-进值特征，而非复值特征。沿用 \ref{VI:1-8}(ii) 的记号，估计 \eqref{VI:eq:7-1-3} 可由下述定理推出。

\begin{theorem_}\label{VI:7-4}
$V_a^{n-1}$ 上以 $\sF(\psi(\sum x_i))$ 为系数的上同调满足：
\begin{enumerate}[\indent (i)]
\item 当 $i\ne n-1$ 时，$\h_c^i=0$；
\item $\h_c^\bullet\iso\h^\bullet$；
\item 若 $a\ne0$，则 $\dim\h_c^{n-1}=n$；
\item 若 $a=0$，则 $\h_c^{n-1}$ 典范同构于 $E_\lambda$。
\end{enumerate}
\end{theorem_}

\subsection{}\label{VI:7-5}
与通常一样，此定理也控制 $\kloos_{n,a}$ 对 $q$ 的依赖：对于每个 $a\in\dF_q^\times$，存在 $n$ 个 Frobenius 特征值 $\alpha_1,\ldots,\alpha_n$，其复绝对值均为 $q^{(n-1)/2}$，使得
\[
\sum_{\substack{x_1\cdots x_n=a\\x_i\in\dF_{q^m}}}
 \psi\circ\tr\left(\sum x_i\right)
=(-1)^{n-1}(\alpha_1^m+\cdots+\alpha_n^m).
\]
当 $n$ 为偶数时，$x\mapsto-x$ 是 $V_a$ 的一个对合，将 $\sF(\psi)$ 变为其对偶。采用与 \ref{VI:3-6} 相同的论证，可在 $\h_c^{n-1}(V_a,\sF(\psi))$ 上得到一个取值于 $E_\lambda(-(n-1))$ 的典范非退化交替型，并得到下述推论。

\begin{corollary_}\label{VI:7-6}
沿用 \ref{VI:7-5} 的记号，若 $n$ 为偶数，则各 $\alpha_i$ 可分成 $n/2$ 对，每对为 $\alpha,\bar\alpha$，乘积等于 $q^{n-1}$。
\end{corollary_}

\subsection{}\label{VI:7-7}
先验证 \ref{VI:7-4} 在 $a=0$ 时的结论。超曲面 $V_0$ 是各坐标超平面 $x_i=0$ 的并。对于以 $\sF(\psi)$ 为系数、有支撑或无支撑的上同调，计算此覆盖的 Leray 谱序列 \eqref{VI:eq:2-6-1*}。由 \ref{VI:2-7*}，初始项全部为零，唯独所有坐标超平面的交集 $\{0\}$ 的上同调例外。这里保留着同构 $\h_c^0\iso\h^0=E_\lambda$，这就证明了 \ref{VI:7-4} 的相应结论。

下面将对 $n$ 作同时归纳，证明 \ref{VI:7-4} 与 \ref{VI:7-8}，所用谱序列正是 \eqref{VI:eq:7-1-4} 的上同调对应物。因此需要控制 $\kloos_{n,a}$ 对 $a$ 的依赖。

设 $\pi:\dA^n\to\dA^1$ 为“各坐标的乘积”映射，$\sigma$ 为各坐标之和。各超曲面 $V_a$ 就是 $\pi$ 的纤维。

\begin{theorem_}\label{VI:7-8}
\begin{enumerate}[(i)]
\item 层 $\eR^{n-1}\pi_!\sF(\psi\sigma)$ 在 $\dA^1\setminus\{0\}$ 上光滑，秩为 $n$；
\item 将此层从 $\dA^1$ 以零延拓到 $\dP^1$，所得层等于它在 $\dA^1\setminus\{0\}$ 上的限制的直像；
\item 在 $0$ 处，单值作用是\term{幺幂的}{unipotent}{unipotent}，且只有一个 Jordan 块；
\item 在 $\infty$ 处，野惯性群的作用没有非零不动向量，且 Swan 导子等于 $1$；
\item $\eR^i\pi_!\sF(\psi\sigma)\iso\eR^i\pi_\ast\sF(\psi\sigma)$，且当 $i\ne n-1$ 时二者均为零。
\end{enumerate}
\end{theorem_}
\begin{proof}
先证明：对于固定的 $n$，\ref{VI:7-4} 蕴涵同一 $n$ 下的 \ref{VI:7-8}。

$\eR^i\pi_!\sF(\psi\sigma)$ 在 $a$ 处的茎为 $\h_c^i(V_a,\sF(\psi\sigma))$。在 $\dA^1$ 的一个稠密开集上，各 $\eR^i\pi_!\sF(\psi^{-1}\sigma)$ 局部常值，而 $\eR^i\pi_\ast\sF(\psi\sigma)$ 的茎在那里为 $\h^i(V_a,\sF(\psi\sigma))$（\hyperref[VII]{Th. finitude}，\ref{VII:2-1}）。假设 \ref{VI:7-4}($n$) 因而给出下述引理。

\begin{lemma_}\label{VI:7-9}
\begin{enumerate}[(i)]
\item 当 $i\ne n-1$ 时，$\eR^i\pi_!\sF(\psi\sigma)=0$；
\item 当 $i=n-1$ 时，该层在每个 $a\ne0$ 处的茎的秩恒为 $n$，在 $0$ 处的秩为 $1$；
\item 在某个稠密开集 $U$ 上，有
$\eR^i\pi_!\sF(\psi\sigma)\iso\eR^i\pi_\ast\sF(\psi\sigma)$。
\end{enumerate}
\end{lemma_}

\subsection{}\label{VI:7-10}
现在使用 $\pi$ 的、有支撑或无支撑上同调的 Leray 谱序列，系数为 $\sF(\psi\sigma)$。其收敛目标已知：
$\h_c^\bullet(\dA^n,\sF(\psi\sigma))=\h^\bullet(\dA^n,\sF(\psi\sigma))=0$
（\ref{VI:2-7*}）。对固有支撑上同调，\ref{VI:7-9}(i) 保证当 $q\ne n-1$ 时 $E_2^{pq}=0$。因此所有 $E_2$ 项均为零：
\begin{equation*}\tag{7.10.1}\label{VI:eq:7-10-1}
\h_c^p(\dA^1,\eR^{n-1}\pi_!\sF(\psi\sigma))=0.
\end{equation*}
当 $p=0$ 时，这意味着 $\eR^{n-1}\pi_!\sF(\psi\sigma)$ 没有点支撑截面。由 \ref{VI:7-9}(ii) 的秩恒定性，得到
\begin{equation*}\tag{7.10.2}\label{VI:eq:7-10-2}
\begin{gathered}
\dA^1\text{ 上的层 }\eR^{n-1}\pi_!\sF(\psi\sigma)
\text{ 在 }\dA^1\setminus\{0\}\text{ 上光滑，}\\
\text{且为其在 }\dA^1\setminus\{0\}\text{ 上的限制的直像 }\sG\text{ 的子层。}
\end{gathered}
\end{equation*}
考虑正合列
\[\xymatrix{0\ar[r]&\eR^{n-1}\pi_!\sF(\psi\sigma)\ar[r]&\sG\ar[r]&\sQ\ar[r]&0,}\]
其中 $\sQ$ 的支撑集中在 $0$。它的上同调长正合列给出
\[\xymatrix{\h_c^0(\dA^1,\sG)\ar[r]&\sQ_0\ar[r]&\h_c^1(\dA^1,\eR^{n-1}\pi_!\sF(\psi\sigma)).}\]
两端均为零，故 $\sQ=0$，于是
\begin{equation*}\tag{7.10.3}\label{VI:eq:7-10-3}
\begin{gathered}
\eR^{n-1}\pi_!\sF(\psi\sigma)\text{ 是其在 }\dA^1\setminus\{0\}
\text{ 上的限制的直像。}
\end{gathered}
\end{equation*}
设 $j:\dA^1\hookrightarrow\dP^1$ 为包含映射。Leray 谱序列的“导出范畴”形式为
\begin{align*}
\hh^\bullet(\dP^1,j_!\eR\pi_!\sF(\psi\sigma))
&=\hh_c^\bullet(\dA^1,\eR\pi_!\sF(\psi\sigma))
=\h_c^\bullet(\dA^n,\sF(\psi\sigma))=0,\\
\hh^\bullet(\dP^1,\eR j_\ast\eR\pi_\ast\sF(\psi\sigma))
&=\hh^\bullet(\dA^1,\eR\pi_\ast\sF(\psi\sigma))
=\h^\bullet(\dA^n,\sF(\psi\sigma))=0.
\end{align*}
令 $\Delta$ 为态射
$j_!\eR\pi_!\sF(\psi\sigma)\to\eR j_\ast\eR\pi_\ast\sF(\psi\sigma)$
的映射锥。由 \ref{VI:7-9}(iii)，这个层复形的上同调层具有有限支撑。另一方面，三角
$(j_!\eR\pi_!\sF(\psi\sigma),\eR j_\ast\eR\pi_\ast\sF(\psi\sigma),\Delta)$
的上同调长正合列表明 $\hh^\bullet(\dP^1,\Delta)=0$。
又有 $\hh^\bullet(\dP^1,\Delta)=\h^0(\dP^1,\sH^\bullet(\Delta))$，最终得到 $\Delta=0$，即
\[j_!\eR\pi_!\sF(\psi\sigma)\iso\eR j_\ast\eR\pi_\ast\sF(\psi\sigma).\]
特别地，$\eR\pi_!\sF(\psi\sigma)\iso\eR\pi_\ast\sF(\psi\sigma)$。由于这些层复形只有一个非零上同调层，故
\begin{equation*}\tag{7.10.4}\label{VI:eq:7-10-4}
j_!\eR\pi_!\sF(\psi\sigma)\iso j_\ast\eR\pi_!\sF(\psi\sigma).
\end{equation*}
这就完成了 (i)、(ii)、(v) 的证明。

\subsection{}\label{VI:7-11}
记 $\swan_0$、$\swan_\infty$ 为 $\eR^{n-1}\pi_!\sF(\psi\sigma)$ 在 $0$ 和 $\infty$ 处的 Swan 导子。由 \eqref{VI:eq:7-10-1}，其 $\chi_c$ 等于零，Euler--Poincaré 公式 \eqref{VI:eq:3-2-1} 因而化为
\begin{equation*}\tag{7.11.1}\label{VI:eq:7-11-1}
\swan_0+\swan_\infty=1.
\end{equation*}
其中一个导子为 $0$，即温和分歧，另一个为 $1$。

\subsection{}\label{VI:7-12}
设 $\sG$ 是 $\dG_m$ 上秩一的非常值 Kummer 型层（\ref{VI:4-7}），即 $\sG=\sK_m(\chi)$，$\chi\ne1$。$V^\ast=\pi^{-1}(\dG_m)=\dA^n\setminus\pi^{-1}(0)$ 上的层
$\sF(\psi\sigma)\otimes\pi^\ast\sG$
属于研究 Gauss 和时遇到的类型，\ref{VI:4-11}(ii) 给出
\begin{equation*}\tag{7.12.1}\label{VI:eq:7-12-1}
\h_c^\bullet(V^\ast,\pi^\ast\sG\otimes\sF(\psi\sigma))
\iso\h^\bullet(V^\ast,\pi^\ast\sG\otimes\sF(\psi\sigma)).
\end{equation*}
$\pi$ 的 Leray 谱序列将此陈述化为
\begin{equation*}\tag{7.12.2}\label{VI:eq:7-12-2}
\h_c^\bullet(\dG_m,\sG\otimes\eR^{n-1}\pi_!\sF(\psi\sigma))
\iso\h^\bullet(\dG_m,\sG\otimes\eR^{n-1}\pi_\ast\sF(\psi\sigma)).
\end{equation*}
设 $i:\dG_m\hookrightarrow\dP^1$，并令 $\Delta$ 为态射
\[\delta:i_!(\sG\otimes\eR^{n-1}\pi_!\sF(\psi\sigma))
\longrightarrow\eR i_\ast(\sG\otimes\eR^{n-1}\pi_!\sF(\psi\sigma))\]
的映射锥，这里 $\eR^{n-1}\pi_!=\eR^{n-1}\pi_\ast$。采用与 \ref{VI:7-10} 相同的论证，可由 \eqref{VI:eq:7-12-2} 得到 $\Delta=0$。特别地，
\begin{equation*}\tag{7.12.3}\label{VI:eq:7-12-3}
\begin{gathered}
\text{在 }0\text{ 和 }\infty\text{ 处，惯性群在 }
\sG\otimes\eR^{n-1}\pi_!\sF(\psi\sigma)\text{ 上的作用}\\
\text{没有非零不动向量。}
\end{gathered}
\end{equation*}
在 $\infty$ 处，即使 $\sG$ 平凡，此结论仍由 \eqref{VI:eq:7-10-4} 成立；因此野惯性群的作用没有非零不动向量。特别地，该处必为野分歧，由 \ref{VI:7-11} 得 $\swan_\infty=1$，而 $0$ 处为温和分歧。

在 $0$ 处，\eqref{VI:eq:7-12-3} 迫使 $\eR^{n-1}\pi_!\sF(\psi\sigma)$ 的温和单值作用为幺幂。由 \eqref{VI:eq:7-10-3} 及 \ref{VI:7-9}(ii)，它只有一个 Jordan 块。至此证明了 \ref{VI:7-8}($n$)。

\subsection{}\label{VI:7-13}
当 $n=1$ 时，\ref{VI:7-4} 显然成立。要完成 \ref{VI:7-4} 和 \ref{VI:7-8} 的证明，还需证明固定 $n$ 时的 \ref{VI:7-8} 蕴涵 $n+1$ 时的 \ref{VI:7-4}。由于 $a=0$ 已在 \ref{VI:7-7} 处理，现假设 $a\ne0$。将 $\dA^{n+1}$ 的坐标记为 $x_0,\ldots,x_n$。设 $g=x_0:V_a\to\dG_m$，$\tau$ 为 $\dG_m$ 的对合 $x\mapsto ax^{-1}$。仍以 $\pi$ 表示坐标乘积映射 $\dA^{n+1}\to\dA^1$，以 $\sigma$ 表示坐标之和。在 $\dA^{n+1}$ 及 $\dA^n$ 上，均将 $\sF(\psi(\sum x_i))$ 记作 $\sF(\psi\sigma)$。

\subsection{}\label{VI:7-14}
写出 $g$ 的 Leray 谱序列。$\dA^{n+1}$ 上的 $\sF(\psi\sigma)$ 是 $\dA^1$ 上的 $\sF(\psi)$ 与 $\dA^n$ 上的 $\sF(\psi\sigma)$ 的外张量积，因此得到（参见 \eqref{VI:eq:7-1-4}）
\begin{align*}
{}'E_2^{pq}
&=\h_c^p(\dG_m,\sF(\psi)\otimes\tau^\ast\eR^q\pi_!\sF(\psi\sigma))
\Rightarrow\h_c^{p+q}(V_a,\sF(\psi\sigma)),\\
{}''E_2^{pq}
&=\h^p(\dG_m,\sF(\psi)\otimes\tau^\ast\eR^q\pi_\ast\sF(\psi\sigma))
\Rightarrow\h^{p+q}(V_a,\sF(\psi\sigma)).
\end{align*}
由假设，$\eR^q\pi_!\sF(\psi\sigma)\iso\eR^q\pi_\ast\sF(\psi\sigma)$，且当 $q\ne n-1$ 时此层为零。$\dG_m$ 上的层 $\sF(\psi)$ 在 $\infty$ 处野分歧、在 $0$ 处不分歧；而 $\tau^\ast\eR^q\pi_!\sF(\psi\sigma)$ 在 $0$ 处野分歧，野惯性群下没有不变向量，Swan 导子为 $1$，在 $\infty$ 处则温和分歧。因此它们的张量积满足：
\begin{enumerate}[\indent a)]
\item 在 $0$ 和 $\infty$ 处均野分歧，且野惯性群下没有不变向量；
\item 在 $0$ 处的 Swan 导子为 $1$，在 $\infty$ 处为 $n$，后者等于 $\eR^{n-1}\pi_!\sF(\psi\sigma)$ 的秩。
\end{enumerate}
故 ${}'E_2^{pq}\iso{}''E_2^{pq}$，且除了 $p=1$、$q=n-1$ 之外，${}'E_2^{pq}=0$。最后，Euler--Poincaré 公式给出
\[\dim{}'E_2^{1,n-1}=1+n.\]
证明完成。
\end{proof}

\subsection{备注}\label{VI:7-15}
\ref{VI:7-14} 所得同构
\begin{equation*}\tag{7.15.1}\label{VI:eq:7-15-1}
\h_c^n(V_a^n,\sF(\psi\sigma))
=\h_c^1(\dG_m,\sF(\psi)\otimes\tau^\ast\eR^{n-1}\pi_!\sF(\psi\sigma))
\end{equation*}
使我们能够沿用 \ref{VI:7-5} 的记号，计算 Frobenius 特征值 $\alpha_i$ 的乘积 $\det(F)$。由于 $\sF(\psi)$ 和 $\eR^{n-1}\pi_!\sF(\psi\sigma)$ 属于无限相容的 $\ell$-进表示系统，\cite{de73} 的形式体系适用。这将问题化为局部问题，而下面两点又可进一步简化它们（参见 \cite[9.5]{de73}）：
\begin{enumerate}[\indent a)]
\item $\sF(\psi)$ 和 $\tau^\ast\eR^{n-1}\pi_!\sF(\psi\sigma)$ 的整体常数容易计算，因为这些层的上同调具有平凡的形式；
\item 在每一点处，这两个层至少有一个所对应的\term{分解群}{decomposition group}{groupe de décomposition} $\ell$-进表示，其\term{半单化}{semisimplification}{semi-simplifiée}不分歧。
\end{enumerate}
对 $n$ 元和，结果为
\begin{equation*}\tag{7.15.2}\label{VI:eq:7-15-2}
\det\left(F,\h_c^{n-1}(V_a^{n-1},\sF(\psi))\right)=q^{n(n-1)/2}.
\end{equation*}

\subsection{备注}\label{VI:7-16}
当 $p=2$ 时，层 $\sF(\psi)$ 是正交的。若 $a\ne0$ 且 $n$ 为奇数，由 \ref{VI:7-4}(ii) 及 Poincaré 对偶，可在 $\h_c^{n-1}(V_a^{n-1},\sF(\psi))$ 上得到取值于 $\bar\dQ_\ell(1-n)$ 的对称双线性型。关于此型，$F$ 是乘子为 $q^{n-1}$ 的正交相似变换；由 \eqref{VI:eq:7-15-2}，$q^{(1-n)/2}F$ 属于特殊正交群。因此，$F$ 有一个特征值为 $q^{(n-1)/2}$，其余特征值可配成 $\alpha,\bar\alpha$ 对，每对乘积为 $q^{n-1}$。
\translatornote{原书在“$F$ 有一个特征值”处误写为 $q^{(1-n)/2}$；此处校为 $q^{(n-1)/2}$。令 $A=q^{(1-n)/2}F$，则 $A$ 是奇数维特殊正交变换，其特征值按倒数配对，剩下的特征值必为 $1$，所以 $F$ 相应的特征值为 $q^{(n-1)/2}$。这也与 \ref{VI:7-5} 中所有特征值的绝对值相符。}
对于 $n=3$，L. Carlitz \cite{ca69-2} 得到了更精确的结果，它等价于下述命题。

\begin{proposition_}\label{VI:7-17}
当 $p=2$ 且 $\psi(x)=(-1)^{\tr_{\dF_q/\dF_2}(x)}$ 时，有同构
\[\h_c^2(V_a^2,\sF(\psi\sigma))
=\operatorname{Sym}^2\h_c^1(V_a^1,\sF(\psi\sigma)).\]
\end{proposition_}
根据 Künneth 公式，$\h_c^1(V_a^1,\sF(\psi\sigma))$ 的二次张量幂为
$\h_c^2(V_a^1\times V_a^1,\sF(\psi\sigma)\boxtimes\sF(\psi\sigma))$。
对称变换 $x\otimes y\mapsto y\otimes x$ 的几何表示是 $-\tau^\ast$，其中 $\tau$ 为 $V_a^1\times V_a^1$ 的自同构 $(x,y)\mapsto(y,x)$。若 $X'$ 为 $V_a^1\times V_a^1$ 对 $\{\operatorname{id},\tau\}$ 的商，$\pi:V_a^1\times V_a^1\to X'$ 为投影，则
$\operatorname{Sym}^2\h_c^1(V_a^1,\sF(\psi\sigma))$
等同于 $X'$ 上以
$\pi_\ast(\sF(\psi\sigma)\boxtimes\sF(\psi\sigma))$
的 $\tau$-反不变部分为系数的上同调。

以 $s$ 表示 $X'$ 上满足 $s\pi=\sigma\pr_1+\sigma\pr_2$ 的函数。有
\begin{align*}
\sF(\psi\sigma)\boxtimes\sF(\psi\sigma)
&=\sF(\psi(\sigma\pr_1+\sigma\pr_2))\\
&=\sF(\psi s\pi)=\pi^\ast\sF(\psi s),
\end{align*}
而且此同构与 $\tau$ 相容。因此，在对角线的像 $\pi(\Delta)$ 上，所求反不变层为零；在补集 $X=X'\setminus\pi(\Delta)$ 上，它等于 $\sF(\psi s)$ 与 $\varepsilon(P)$ 的张量积，其中 $\varepsilon$ 是 $\{1,\tau\}$ 的唯一二阶特征，$P$ 是由 $V_a^1\times V_a^1$ 诱导的二重覆盖，即 $X$ 上的一个 $\{1,\tau\}$-扭子。于是
\begin{equation*}\tag{7.17.1}\label{VI:eq:7-17-1}
\operatorname{Sym}^2\h_c^1(V_a^1,\sF(\psi\sigma))
=\h_c^2(X,\sF(\psi s)\otimes\varepsilon(P)).
\end{equation*}
下面计算 $X,s,\varepsilon(P)$。在 $V_a^1\times V_a^1$ 上，令
$x_1=x_1\circ\pr_1$、$x_2=x_2\circ\pr_1$、$x_1'=x_1\circ\pr_2$、$x_2'=x_2\circ\pr_2$。有 $x_1x_2=x_1'x_2'=a$。将 $X'$ 上的函数等同于 $V_a^1\times V_a^1$ 上的 $\tau$-不变函数。令
$s_1=x_1+x_1'$、$s_2=x_2+x_2'$、$p_1=x_1x_1'$、$p_2=x_2x_2'$。则
$s_1p_2=as_2$、$s_2p_1=as_1$，而 $X$ 在 $X'$ 中由 $s_1\ne0$、$s_2\ne0$ 定义。因此，在 $X$ 上有 $p_1=as_1s_2^{-1}$，$(s_1,s_2)$ 将 $X$ 等同于 $\dG_m\times\dG_m$。在 $X$ 上，覆盖 $P$ 的方程为
$T^2-s_1T+as_1s_2^{-1}=0$，其中取 $T=x_1$。作代换 $T=s_1t$，方程成为
$t^2-t+as_1^{-1}s_2^{-1}=0$。因而
$\varepsilon(P)=\sF(\psi(as_1^{-1}s_2^{-1}))$。
由于 $s=s_1+s_2$，有
\[\sF(\psi s)\otimes\varepsilon(P)
=\sF(\psi(s_1+s_2+as_1^{-1}s_2^{-1})).\]
因此 \eqref{VI:eq:7-17-1} 的右侧等同于命题 \ref{VI:7-17} 的左侧。

\paragraph{备注}
$\h_c^1(V_a^1,\sF(\psi\sigma))$ 也是仿射曲线
\begin{align*}
y^2-y&=x_1+x_2,\\x_1x_2&=a
\end{align*}
的光滑射影完备化所得椭圆曲线的一次上同调群；其模不变量为 $j=a^{-2}$。

\subsection{备注}\label{VI:7-18}
设 $\chi$ 为 $\dF_q^\times$ 的乘法特征，$a\ne0$。\ref{VI:7-14} 的方法也适用于研究
\[s=\sum_{x_0\cdots x_n=a}\chi(x_0)\psi(x_0+\cdots+x_n),\]
只需将 $\dG_m$ 上的 $\sF(\psi)$ 换成 $\sF(\chi\psi)$。仍得到维数为 $n+1$ 的上同调，且 $|s|\le(n+1)q^{n/2}$。Salié 和 Mordell 曾考虑过这样的和 \cite{mo73}。

可以期望，我们的方法也能用于研究
\[S_{k,\chi,a}=\sum_{N(x)=a}\chi(x)\psi(\tr(x)),\]
其中 $k$ 为 $\dF_q$ 上的平展代数，$a\in\dF_q^\times$，$\chi$ 为 $k^\times$ 的特征。\ref{VI:7-4} 在 $a\ne0$ 情形的对应结论应当成立。

\subsection{}\label{VI:7-19}
对称群 $S_n$ 在 $\dA^n$ 上的作用保持 $V_a$ 以及映射 $\sigma$，因而也保持层 $\sF(\psi\sigma)$，即 $\dG_a$ 上的 $\sF(\psi)$ 沿 $\sigma$ 的逆像层。因此，它在 $\h_c^{n-1}(V_a,\sF(\psi\sigma))$ 上作用。公式 \eqref{VI:eq:7-2-5} 具有如下上同调解释。

\begin{proposition_}\label{VI:7-20}
$S_n$ 在 $\h_c^{n-1}(V_a,\sF(\psi\sigma))$ 上的作用是乘以符号特征。
\end{proposition_}
只需证明换位 $\tau$ 的作用是乘以 $-1$。由于 $\tau^2=1$，$\tau$ 的特征值只能为 $\pm1$。设 $V_a/\tau$ 为 $V_a$ 对 $\{\operatorname{id},\tau\}$ 的商。映射 $\sigma$ 以及层 $\sF(\psi\sigma)$ 均下降到此商，$\h_c^{n-1}(V_a,\sF(\psi\sigma))$ 的 $\tau$-不动子空间为
$\h_c^{n-1}(V_a/\tau,\sF(\psi\sigma))$。只需证明该上同调为零。

设 $a\ne0$。对于 $\tau=(1,2)$，商 $V_a/\tau$ 等同于 $\dG_a\times\dG_m^{n-2}$，取商映射为
$(x_1,x_2,x_3,\ldots,x_n)\mapsto(x_1+x_2,x_3,\ldots,x_n)$。
事实上，$x_1+x_2$ 和 $x_1x_2=a/(x_3\cdots x_n)$ 决定了 $x_1,x_2$，只差次序。在这些坐标 $(t,x_3,\ldots,x_n)$ 下，层写成
$\sF(\psi(t+x_3+\cdots+x_n))=\sF(\psi t)\otimes\sF(\psi(x_3+\cdots+x_n))$。
由 Künneth 公式及 $\h_c^\bullet(\dG_a,\sF(\psi))=0$，即得所需的消失结论。

当 $a=0$ 时，可由 \ref{VI:7-7} 的计算推出 \ref{VI:7-20}。

\subsection{}\label{VI:7-21}
S. Sperber 在其学位论文以及即将发表的 \emph{p-adic hypergeometric functions I, II} 中，采用受 Dwork 及 Bombieri 的论文 \emph{On exponential sums in finite fields}（\emph{Amer. J. Math.} 88，1966）启发的 $p$-进方法，研究了本节的和 $\kloos_{n,a}$。当 $p\ne2$ 时，这些方法给出了 \ref{VI:7-5} 中对 $q$ 的依赖关系、当 $b=(-1)^na$ 时和 $\kloos_{n,a}$ 与 $\kloos_{n,b}$ 的 Frobenius 特征值之间的函数方程
$\alpha_i=q^{n-1}\beta_i^{-1}$（参见 \ref{VI:7-6}），以及 \ref{VI:7-15} 中各根的乘积。当 $p\ge n+3$ 时，他还确定了各根的 $p$-进赋值。结果非常漂亮：适当排列各根 $\alpha_i$，以 $0\le i<n$ 编号后，各 $\alpha_iq^{-i}$ 均为单位。这里对 $p$ 的限制大概不是本质的。

\section{其他应用}\label{VI:8}
\origpage{236}{230}

\subsection{}\label{VI:8-1}
恒等式 \ref{VI:1-9-4} 有时可以用来计算 $\dF_q$ 上代数簇的有理点个数。在大多数已经获得精确结果的情形中，上同调都可以用代数闭链表达；有理曲面、二次超曲面、奇数维的两个二次超曲面的光滑交、旗簇等都属此类。不过这并非总是如此；例如在 Lusztig 那篇非常漂亮的论文 \cite{lu76} 中，代数闭链并未出现。

\subsection{}\label{VI:8-2}
对于\term{约化群}{reductive group}{groupe réductif}，可以通过提升到特征 $0$ 来计算上同调。若 $B$ 是 $G$ 的一个 Borel 子群，则 $G$ 是 $G/B$ 上的一个 $B$-扭子。光滑射影簇 $G/B$ 的上同调在特殊化下不变，再用 $G\to G/B$ 的 Leray 谱序列便可对 $G$ 证明相同的结论。此上同调可方便地用 $G$ 的那个\term{极大环面}{maximal torus}{tore maximal} $T$ 以及那个 \term{Weyl 群}{Weyl group}{groupe de Weyl} $W$ 在 $T$ 上的作用来表达（关于“那个”极大环面的定义，见 \cite[VIII，\S2，备注 2]{bo05} 或 \cite[p.105]{dl76}）。它是其本原部分的外代数，而该本原部分在差一个平移的意义下，等于 $BG$ 的上同调中由“不可分解元”构成的商
$I(\h^\bullet(BG,\dQ_\ell))$；又有
$\h^\bullet(BG,\dQ_\ell)=\h^\bullet(BT,\dQ_\ell)^W$。
若 $X$ 是 $T$ 的特征群，综上得到
\[
\h^\bullet(G,\dQ_\ell)=\bigwedge\left(
 I\left(\operatorname{Sym}^\bullet(X\otimes\dQ_\ell(-1))^W\right)[-1]\right).
\]
若 $G$ 在有限域 $\dF_q$ 上有定义，则 $\gal(\dF/\dF_q)$ 在 $X$ 上作用，且此同构与 Galois 作用相容。由此转到固有支撑上同调，得到 $G$ 的有理点个数的经典公式。若
$F^\ast:X\to X$ 由 Frobenius 态射 $F:T\to T$ 定义，且 $\deg_G(F)$ 为 $F:G\to G$ 的次数 $q^{\dim G}$，则
\begin{equation*}\tag{8.2.1}\label{VI:eq:8-2-1}
|G_0(\dF_q)|=\deg_G(F)\cdot
\det\left(1-(F^\ast)^{-1},I(\operatorname{Sym}^\bullet(X\otimes\dQ)^W)\right).
\end{equation*}
Ree 群和 Suzuki 群也适用同一公式：它们具有 $G^F$ 的形式，其中 $G$ 为约化群，$F^2$ 为一个 Frobenius。只需验证 $F$ 在 $G$ 上作用时 Lefschetz 迹公式成立。若 $B$ 是 $F$ 稳定的 Borel 子群，则 $F$ 在 $B$ 上的作用可直接验证，而在固有光滑的 $G/B$ 上的作用则可使用一般定理。

\subsection{}\label{VI:8-3}
Deligne--Lusztig \cite{dl76}、Kazhdan \cite{ka77} 和 Springer \cite{sp76} 已将迹公式用于研究有限群 $G_0(\dF_q)$ 的复表示，其中 $G_0$ 是 $\dF_q$ 上的约化群。Kazhdan 和 Springer 的工作提供了一些极为精彩的例子，展示如何用迹公式将分裂情形“解析延拓”到非分裂情形（参见 \ref{VI:1-13}）。

\section*{参考文献}
\origpage{238}{232}
\begin{description}
\item[{[1]}] L. Carlitz，\emph{A note on exponential sums}，\emph{Pacific J. Math.} 30 (1969)，35--37 页；MR 39 6839。
\item[{[2]}] P. Deligne，\emph{Les constantes des équations fonctionnelles des fonctions L}，Proc. Antwerpen conf.，第 2 卷，501--597 页，Lecture Notes 349，Springer-Verlag。
\item[{[3]}] P. Deligne、G. Lusztig，\emph{Representations of reductive groups over finite fields}，\emph{Ann. of Math.} 103，1 (1976)，103--161 页。
\item[{[4]}] D. Kazhdan，\emph{Proof of Springer's Hypothesis}，Harvard 预印本。
\item[{[5]}] J. P. Serre，\emph{Abelian $\ell$-adic representations and elliptic curves}，Benjamin，1968。
\item[{[6]}] T. A. Springer，\emph{Trigonometrical sums, Green functions and representations of Weil groups}，Utrecht 预印本。
\item[{[7]}] A. Weil，\emph{Basic number theory}，Springer-Verlag。
\item[{[8]}] A. Weil，\emph{Jacobi sums as ``Grössencharaktere''}，\emph{Trans. Amer. Math. Soc.} 73 (1952)，487--495 页。
\item[{[9]}] A. Weil，\emph{Sommes de Jacobi et caractères de Hecke}。
\end{description}
